% 论文主控 tex（用户直接在此手改）；上游 md2latex 生成流程已停用，勿再重新生成覆盖
\documentclass[12pt,a4paper]{ctexart}
\usepackage{geometry}
\geometry{left=2.8cm,right=2.8cm,top=2.5cm,bottom=2.5cm}
\usepackage{booktabs}
\usepackage{tabularx}
\usepackage{array}
\usepackage{amsmath,amssymb}
\setlength{\parindent}{2em}
\setlength{\parskip}{4pt}
\linespread{1.25}
\setcounter{secnumdepth}{0}
\setcounter{tocdepth}{2}
\renewcommand{\arraystretch}{1.25}
\begin{document}
\begin{titlepage}
\vspace*{3cm}
\begin{center}
{\LARGE\bfseries 数据形状与危机识别——拓扑数据分析的数学原理、实验证据与产品化实践\par}
\vspace{2.2cm}
{\normalsize \textbf{作者}：潘炯铠\par\vspace{0.35em}}
{\normalsize \textbf{单位}：持续同调量化研究项目 · 工作论文\par\vspace{0.35em}}
{\normalsize \textbf{数据区间}：原始日历 1999-12-30 至 2026-11-03（6,502 个交易日）；训练期为 2001-01-10 至 2025-04-07（45 个按固定标签规则识别的危机事件，末次事件 2025-03-28 结束、其标签延伸至 2025-04-07）；2025-04-08 起为验证期（样本口径详见引言）\par\vspace{0.35em}}
{\normalsize \textbf{成文日期}：2026-11-16\par\vspace{0.35em}}
\vspace{1.6cm}
\begin{minipage}{0.86\textwidth}\small\setlength{\parindent}{2em}
\textbf{声明}：本文为工作论文，未经同行评议；全部证据边界与历史重放口径的完整声明见引言后「研究声明」。\par
\end{minipage}
\end{center}
\end{titlepage}
\tableofcontents
\newpage
\begin{abstract}
拓扑数据分析（TDA）试图从市场数据形状的变化中识别危机信号，但现有研究对其预警效果尚无一致结论，也很少有研究回答一个更实际的问题：把拓扑信号用于仓位、止损与对冲决策，扣除误报成本后究竟有没有价值。本文据此构建了一套 A 股危机识别与动态风控系统。系统每天收盘后回答三个问题：是否需要提高警惕、当前市场形态更接近哪一类、应切换到哪一种风控预案。它不追求精确预测崩盘日期，而是将触发告警、状态判别、告警撤销与冷却、损益核算连接为一条可逐项检查的流程。与任何危机事件的规定时间范围均无交集的连续告警段被计为误报；损益评价同时计入危机期间的预案效果与这些误报造成的成本。

在 2001—2025 年按固定规则标记的 45 个 A 股市场危机事件上，冻结基准方案采用「三项价格与波动条件满足其二即告警」的日频规则，并以「连续 5 日无命中即撤销、撤销后冷却 10 个交易日」管理告警状态。统计是否识别到危机时，允许信号最早出现在事件起点前 20 个交易日，最晚出现在事件结束后 20 个交易日。按这一口径，45 次危机均有告警覆盖：24 次在下跌开始前触发，7 次在下跌过程中触发，14 次直到下跌结束后才识别到。因此，全部识别不等于全部提前预警。45 次中有 31 次来得及在事件结束前启动预案核算，实际切换风控预案的为 19 次；告警日占全部交易日的 23.5\%。样本内历史重放的净损益差合计为 +79.0，其中危机事件内的预案损益差为 +147.4，误报段损益差为 \(-\)68.4。该数值是把各事件与各误报段从单位净值出发计算的损益差逐项相加得到的对照结果，\textbf{不是单一组合的累计收益率}。

在候选组合中，方案二与分段重新计时规则的组合取得最高样本内净损益差（+140.7）；由于两项修订均经过样本内选择，且验证期未显示类型判别改善，本文将其定位为后续前瞻验证的首选组合，而不替代冻结基准方案。

方案一曾在事先留出的验证区间（2025-04-08 起）进行封存样本外检验：区间内 6 次危机全部被识别，其中 2 次预警信号领先指数拐点 20—23 个交易日；经过完整的告警撤销与冷却后，非危机区间的年均误报天数从训练期的 29.9 天降至 3.8 天，2026 年全年零误报。

结果表明，TDA 不宜作为独立预测崩盘时点的工具，更适合作为对市场相关结构变化敏感的信息增强模块。后续候选方案再次使用了上述验证区间，因此相关结果只能视为描述性重放，不能算新的独立样本外证据。当前系统仍存在实时类型判别精度不足、历史重放依赖事后已知的事件结束日、未计一般换手摩擦等局限；全部样本内绩效均为上限口径，证据边界在第五章末集中列明。\par
\textbf{关键词}：市场状态识别；持续同调；拓扑数据分析；危机预警；历史重放损益评估；误报控制\par
\end{abstract}
\section{引言}
\subsection{拓扑数据分析在各科学领域的应用}
数据的「形状」蕴含信息，是拓扑学进入数据分析领域的基本信念。拓扑数据分析（下称 TDA）以代数拓扑为数学基础，从高维、含噪数据中提取对坐标表示与连续形变不敏感的结构特征——连通分支、环、空腔及其在多尺度下的生灭规律。由于其输出不依赖人为阈值、对缺失与噪声样本相对稳定，近二十年来 TDA 已从纯数学工具成长为跨学科的标准分析手段。\par
在\textbf{神经科学}领域，马克斯·普朗克数学科学研究所等机构的研究者用持久同调分析大脑功能连接网络，在多尺度下追踪拓扑特征的重构规律，并提出「节点持续性」度量，识别出连接性随年龄变化或在自闭症谱系障碍人群中存在差异的脑区，为无创脑刺激疗法提供了潜在靶点（Patterns，2025）。与依赖任意阈值的传统网络分析相比，持久同调为大脑网络提供了稳健、无参数的特征描述。在\textbf{生物医学}领域，哥伦比亚大学团队于 2017 年在《自然·生物技术》提出单细胞拓扑数据分析算法，重建干细胞分化的发育轨迹；后续研究将 TDA 用于癌症基因组的变异识别（Nature Communications，2020）与脊髓损伤、创伤性脑损伤的临床前数据分析（Nature Communications，2015）。在\textbf{材料科学与凝聚态物理}中，拓扑方法用于理解拓扑相变与拓扑物态，解释宏观性质对局部缺陷的高稳定性。在\textbf{气候与复杂系统}领域，拓扑—动力学共演化的视角被用于研究气候系统振荡与系统重构。在\textbf{网络科学、社交网络、城市结构与图像分析}中，单纯复形被用来建模多体协同行为，揭示传统成对网络看不见的高阶结构（三角形、四面体等），而点云数据的拓扑重建已是三维扫描等场景的基础工具。\par
这些领域的共同经验是：TDA 的价值不在替代传统统计量，而在补足它们看不见的\textbf{高阶结构与多尺度生灭}信息；同时，数学上的稳定性定理保证了拓扑摘要对输入扰动的连续依赖，这是统计推断得以立足的前提。\par
\subsection{量化投资领域：从市场状态到危机识别}
TDA 进入金融领域较晚——文献记载的第一篇金融 TDA 论文出现在 2017 年前后——但发展迅速，应用覆盖市场状态识别、危机预警、波动率与尾部风险预测、组合构建、方向预测等多个方向。其中\textbf{危机识别与危机预测}是最受关注、也最富争议的方向。\par
范式工作由 Gidea 与 Katz 奠定：对股指日对数收益序列做滑动窗口延迟嵌入构造点云，计算持续同调并以持久景观的范数构造时间序列，发现 2000 年互联网泡沫破裂与 2008 年雷曼破产前，景观范数持续走高；其中刻画范数长期趋势的低频成分，在两次崩盘前约 250 个交易日持续上升，并在崩盘中达到峰值（Physica A，2018）。后续扩展包括：Guo 等将分析推广到美欧中多市场并比较多个同调维度；Aguilar 与 Ensor 引入均值持久景观与置换检验；Ismail 等将景观范数与临界慢化指标结合，并同时评估「成功预警概率」与「错误预警概率」，诚实暴露了误报负担——同团队的后续研究（Frontiers in Applied Mathematics and Statistics，2022）报告：即使在最佳设定下，错误预警的概率仍超过一半；Baitinger 与 Flegel 以相邻持久图之间的威尔斯特拉斯距离构造拓扑湍流指数，并做样本外状态预测检验，声称优于其他流行的湍流指数（Financial Markets and Portfolio Management，2021）；Guritanu 等给出含因果决策规则、基线对比与消融设计的检测协议模板，是该方向方法学上最完整的工作（Computers，2025）。针对极端事件，Gidea 等对 2010 年美股闪崩的专项研究、以及结合小波变换自动选取嵌入窗口的泡沫检测新作（2025 年金融数学会议），代表了该路线的最新推进。\par
然而，\textbf{反面证据同样扎实}。Valdivia 在 2018、2020、2022 三个事件上未能得到一致的拓扑响应（Quantitative Finance and Economics，2023）；Ruiz-Ortiz 等指出部分指数的拓扑峰值出现在危机开始之后，属于滞后指标，对缓慢下跌尤其迟钝（arXiv，2022）；预警理论文献则警示「只在已知危机上事后验证」的检察官谬误（Boettiger 与 Hastings，2012）——该偏倚会严重高估预警指标的效能。\par
归结起来，文献留下的三个空白是：（一）\textbf{增量}——拓扑摘要在常规风险指标（波动率、相关性、离散度）之上能否提供稳定的危机预测增量，鲜有严格检验；（二）\textbf{异质性}——拓扑响应方向是否随危机类型而变，这可能是文献方向报告互相打架的根源；（三）\textbf{经济价值}——拓扑信号接入组合、计入误报与交易成本后值多少钱，认真算过这笔账的工作很少。\par
\subsection{本文工作}
本文不提出新的拓扑不变量，并非纯数学层面的理论研究。文章围绕四个相关问题展开，但篇幅并非平均分配：核心成果是一套\textbf{危机识别与动态风控系统}——简言之，系统以日度频率回答三个问题：当前市场是否处于危机区间、危机属于哪种类型、对应需要启动什么等级的风控预案；同时通过告警状态管理与冷却规则减少重复激活及其成本，并对每一段告警的收益与成本做可追溯的绩效核算（整体采用「触发识别—状态判别—告警抑制—绩效核算」四层架构，具体实现见第五、六章）。前述关于危机异质性、策略经济价值的两个问题，都由这套主线系统来回答；而关于 TDA 是否能独立提供增量信息的问题，由两条独立研究支线补充说明，这部分内容不会进入后续摘要，在此做明确交代：

第一条支线是\textbf{对 “直接用拓扑特征预测危机” 路线的证伪（第三章）}。我们先搭建基于五项常规风险指标的逻辑回归基线模型，再在模型中加入拓扑结构特征做对比，结果显示特征加入后并未带来稳定的危机预测能力提升：精确率—召回率曲线下面积的增量为 \(-\)0.0076，95\% 分块自助置信区间为 [\(-\)0.0207, +0.0143]。按照研究预先设定的停止规则，我们终止了 “直接拿拓扑摘要预测危机发生” 的研究路线；需要明确的是，这一结果仅说明该特定技术路线没有达到预期效果，并不构成 “拓扑特征完全不包含额外市场信息” 的最终结论。

第二条支线是\textbf{对券商公开研报策略的独立复现与检验（第四章）}。我们复现了兴业证券提出的 “低拓扑孔洞数做多”（即基于二维网络孔洞数 \(N \le 1\) 时抄底）的交易规则：独立计算发现，该规则的日度平均信息系数确实为正（区间 0.09—0.13），但 “拿到规则就能直接交易” 的可交易性结论未通过进一步检验：在更贴合实际交易决策的信号首次触发口径下，规则的信息系数降至 0.013，且置信区间跨零、不再具备统计显著性；原研报给出的固定触发阈值在 2020 年 5 月之后几乎不再触发；且对比常规超跌反弹基准策略，该规则没有提供独立的收益增量。这一结果说明，规则本身确实携带部分市场信息，但将其视为独立且具备跨期稳定性的实盘交易信号并不成立。两条支线的共同价值是为后续的主线系统划清了能力边界：拓扑特征并不适合直接预测崩盘时点，更适合在告警触发后对市场状态做精细判别 —— 这也是我们最终把拓扑特征放在 “状态判别” 模块，而非 “崩盘预测” 模块的核心原因。

本文的研究对象覆盖 A 股全市场：原始交易日历共 6,502 天（1999-12-30 至 2026-11-03），样本包含 16 个早期、逐步扩容至 31 个的申万一级行业面板数据，以及基于行业间相关性构建的网络结构；危机标签采用预先冻结的标准：未来 20 个交易日内最低累计收益 \(\leq\) \(-\)10\% 即标记为危机事件，共识别出 45 个研究事件（时间跨度从 2001-01-10 至 2025-03-28，对应事件标签最远更新至 2025-04-07）。

整个实验严格遵守样本纪律：所有用于模型开发、参数校准的样本均为开发段数据，而非用全部交易日做训练。其中第一条支线的开发期样本外评估区间为 2008-02-14 至 2024-12-03，券商策略复现的检验区间为 2015-07-01 至 2025-05-30；方案一冻结规则的训练与校准截止时间为 2025-04-07（最后一次危机事件及其对应标签的更新时间），从 2025-04-08 起为留出验证期，该区间的封存样本外评估已完成一次，结果见第六章。第五章确认层扩展候选（方案二）的补全实验再次使用该区间，仅属描述性重放，不能再视为未被使用过的独立验证集。其中第一条支线的实验遵守更严格的封存规则：特征、标签、评分对应的最终测试集从 2025-01-01 起完全封存，在最终评估前从未读取评分结果。

我们需要如实说明统计口径：逐日触发、告警状态管理、误报统计与主损益核算均按照预登记规则在全交易日历上开展（包含 2025 年 4 月之后的无事件交易日，统计最远延伸至 2026-11-03），因此“年均误报天数”这类比率的分母包含留出验证期的交易日。验证期公开行情不参与模型训练与参数校准，但会进入误报统计、全历史损益核算及单独披露的验证期重放；主损益结果因此不是纯训练期隔离结果，也不能与验证期重放结果重复相加。整个研究流程坚持“先预登记、后开展计算”，所有实验结果经过两次独立运行做一致性校验，未达预期的负面结果也全部如实纳入正文。

全文章节安排如下：第一章介绍 TDA 方法的核心数学原理；第二章梳理量化投资领域已有 TDA 相关研究的贡献与局限；第三章汇报 “直接用拓扑特征预测危机” 路线的证伪证据及其结论边界（对应第一条支线）；第四章汇报兴业证券公开策略的复现与可交易性检验结果（对应第二条支线）；第五章介绍规则型危机状态判别器的设计与实证结果；第六章介绍基于该系统落地的决策应用与留出样本验证；附录汇总本文核查过的金融领域相关 TDA 文献及其核心结论。

\vspace{0.8em}
\noindent\rule{\textwidth}{0.6pt}
\vspace{0.6em}

\noindent\begin{minipage}{0.96\textwidth}\small\setlength{\parindent}{2em}\sloppy
\textbf{研究声明}\quad 本文为工作论文，未经同行评议。训练期及全历史重放结果均为样本内描述性统计与历史重放损益评估。方案一曾在 2025-04-08 起的预留区间进行过一次封存样本外检验。第五章各候选修订再次使用该区间，结果仅属描述性重放，不构成新的独立样本外证据。本文不构成实盘业绩宣称或投资建议。\par
\end{minipage}

\vspace{0.5em}
\noindent\rule{\textwidth}{0.6pt}

\par
\par
\section{第一章　拓扑数据分析的数学原理}
本章给出拓扑数据分析（下称 TDA）从数据到统计量的完整数学链条：先把点云组织成随尺度生长的单纯复形，用同调群量化各维「洞」，用持久同调记录洞的生灭，再把生灭记录压缩成可统计的摘要（持久图、持久景观），最后由稳定性定理保证这些摘要对噪声不敏感。本章记号与标准教材一致（Edelsbrunner 与 Harer，2010；Chazal 与 Michel，2021），景观部分依 Bubenik（2015），稳定性定理取 Cohen-Steiner 等（2007）的原始形式；后文全部拓扑特征均以此处的定义为准。\par
\subsection{1.1　从数据到单纯复形}
设数据集 X 为度量空间中的有限点集。单纯复形是把「点—边—三角形—四面体……」逐级粘合而成的组合对象：k 维单纯形是 k+1 个仿射无关点的凸包；单纯复形 K 是一族单纯形的集合，对取面封闭（任何单纯形的面仍在 K 中），且任意两个单纯形或不交、或交于公共面。\par
TDA 不直接对 X 做同调，而是先用尺度参数 r 把 X 转成一族随 r 单调生长的复形。两种标准构造：\textbf{切赫（Čech）复形}以每点为心、r/2 为半径作球，把球族有公共交的指标集收作单纯形；神经定理保证它与球族并集同伦等价，是「多尺度重建形状」的理论原型。\textbf{维耶托里斯–里普斯（Vietoris–Rips）复形}则只看两两距离：若 X 的子集中任意两点距离不超过 r，则该子集构成的单纯形全部收入复形，
\[
VR_r(X) = \{\sigma \subseteq X : d(x,y) \le r,\ \forall x,y \in \sigma\}.
\]
本文取里普斯构造：不依赖嵌入、组合上易算，适合行业收益这样的高维点云。随 r 从 0 递增，复形只增不减，形成滤过序列 \(VR_{r_1}(X) \subseteq VR_{r_2}(X)\)（当 \(r_1 \le r_2\)）。直观上 r 是「分辨率」：r 很小时点云只是散点；r 增大，邻近点先连边、再填三角形；r 足够大时一切连成一团。持久同调要量化的，正是这团「生长史」中洞的生灭。\par
\subsection{1.2　同调群与贝蒂数}
固定一个域 \(\mathbb{F}\)（本文实现取 \(\mathbb{F}_2\)，免去挠元的技术细节）。对复形 K，k 维链群 \(C_k(K)\) 是 k 维单纯形的形式和构成的向量空间；边界算子 \(\partial_k : C_k \to C_{k-1}\) 由
\[
\partial_k[v_0,\ldots,v_k] = \sum_{i=0}^{k} (-1)^i [v_0,\ldots,\widehat{v_i},\ldots,v_k]
\]
给出（\(\widehat{v_i}\) 表示删去该顶点），且满足 \(\partial_{k-1} \circ \partial_k = 0\)。由此定义循环 \(Z_k = \ker \partial_k\)、边界 \(B_k = \operatorname{im} \partial_{k+1}\)，以及 k 维同调群 \(H_k(K) = Z_k / B_k\)。\par
贝蒂数 \(\beta_k = \dim_{\mathbb{F}} H_k\) 计数「本质的洞」：\(\beta_0\) 是连通分支个数，\(\beta_1\) 是独立环路（一维洞）个数，\(\beta_2\) 是空腔（二维洞）个数。经典例子：三个顶点用边连成空心三角形，\(\beta_1 = 1\)（圈内有洞）；把三角形内部填实，\(\beta_1 = 0\)——洞的死灭就是被更高维单纯形填上的时刻。市场数据的常用读法是：\(\beta_0\) 反映资产群体的分裂与合并，\(\beta_1\) 反映轮动与循环结构，\(\beta_2\) 反映更高阶的协同空腔。\par
\subsection{1.3　持久同调、持久图与条形码}
单个尺度的 \(\beta_k\) 只是一张快照；滤过给了整条生长曲线。对 \(r_1 \le r_2\)，包含映射 \(VR_{r_1} \hookrightarrow VR_{r_2}\) 诱导同调间的线性映射 \(H_k(VR_{r_1}) \to H_k(VR_{r_2})\)，全体 \(\{H_k(VR_r)\}\) 连同这些映射构成一个持久模。持久同调的区间分解定理说：持久模可唯一分解为若干区间模的直和，每个区间对应一个洞的存活区间 \([b_i, d_i)\)——\(b_i\) 是出生尺度，\(d_i\) 是被填上（死亡）的尺度。全体区间构成条形码（barcode），等价地记作持久图 \(\mathrm{Dgm}_k = \{(b_i, d_i)\}_i\)。\par
持久性定义为 \(d_i - b_i\)：跨尺度持续越久的特征越可能是真实结构，短命特征多为噪声。常用标量摘要包括总持久性 \(A = \sum_i (d_i - b_i)\)、最大持久性、各尺度下的贝蒂数曲线 \(\beta_k(r) = \#\{i : b_i \le r < d_i\}\) 等。需要说明：每个连通分支都在 \(r = 0\) 出生，其中除一个之外都在某尺度并入别的分支而死亡，另有一个随 r 增大永不死亡（本质类）；实践中对本质类单独处理或截断，不影响有限持久性特征的统计。持久同调的可贵之处在于\textbf{不需要人为划分信号与噪声的阈值}——所有特征保留，按持久性加权。\par
\subsection{1.4　持久景观：把持久图搬进函数空间}
持久图是多重集，空间结构生硬：两张图不能直接取平均，均值、方差、主成分回归都无从谈起。Bubenik（2015）的持久景观（persistence landscape）把它变换为函数序列。对每个出生—死亡对 \((b, d)\)，定义帐篷函数
\[
f_{(b,d)}(t) = \max\{0,\ \min(t - b,\ d - t)\},
\]
即区间 \([b, d]\) 上以中点为峰、峰高 \((d - b)/2\) 的三角形。第 k 层景观函数取所有帐篷函数在 t 处的第 k 大值：
\[
\lambda_k(t) = \underset{i}{k\text{-th largest}}\ f_{(b_i, d_i)}(t),\qquad k = 1, 2, \ldots
\]
景观即序列 \(\lambda = (\lambda_1 \ge \lambda_2 \ge \cdots)\)。它落入巴拿赫空间 \(L^p\)（离散化后即有限维向量），于是可以定义景观范数 \(\|\lambda_k\|_p\)、逐点平均景观、期望与方差，并满足大数定律与中心极限定理——这是拓扑摘要进入统计推断的标准接口，也是金融文献（如 Gidea 与 Katz 的「崩盘景观」）直接引用的量。景观还继承稳定性：两张持久图的景观在一致范数下的距离不超过它们的瓶颈距离。\par
\subsection{1.5　距离与稳定性定理}
持久图之间的比较用最优匹配定义。\textbf{瓶颈距离} \(d_B\)：在两张图的点之间做部分匹配，允许多余点匹配到对角线（代价取该点到对角线的距离，即持久性之半），取所有匹配中最大代价的下确界。\textbf{威尔斯特拉斯距离} \(W_p\) 把匹配代价取 p 次幂求和再开 p 次方根；该距离以「把一张持久图的质量搬运到另一张」的最优代价定义，是持久图空间上的度量，本文统一译名「威尔斯特拉斯距离」，并用二阶距离 \(W_2\) 刻画相邻交易日持久图的变化速度。\par
稳定性定理（Cohen-Steiner 等，2007）是 TDA 可用于含噪数据的根本理由。对正规函数 f、g，持久图的瓶颈距离不超过函数的上确界范数距离；对点云版本，以豪斯多夫距离 \(d_H\) 控制：
\[
d_B\big(\mathrm{Dgm}(X),\ \mathrm{Dgm}(Y)\big) \le 2\, d_H(X, Y)\qquad(\text{里普斯滤过}).
\]
切赫滤过下系数为 1；威尔斯特拉斯距离亦有同型上界。换言之，\textbf{输入只动一点，拓扑摘要最多动一点}——采样噪声不会被放大为虚假的拓扑突变。等价的模语言是：\(\delta\)-交织的持久模，其瓶颈距离不超过 \(\delta\)（交织距离与瓶颈距离等距）。\par
必须强调一条经常被误读的界线：\textbf{数学稳定性不等于金融预测的样本外稳定性}。前者是拓扑定理，保证「测量工具不抖」；后者是时间序列性质，必须靠样本外验证、强基线与交易成本证明。本文第三章的负面结果正是这条界线的实证注脚。\par
\subsection{1.6　本文的两类点云构造}
\textbf{市场状态点云}。一点一交易日，坐标为各行业日对数收益向量 \(r_t \in R^{n}\)（n 为行业数，16/27/31）。t 日的市场状态点云取最近 \(W = 50\) 日：\par
\(X_t = \{ r_{t-W+1}, \ldots, r_t \} \subset R^{n}\)。\par
点间取欧氏距离。该构造把「最近 50 日市场每天的横截面形态」视为一个几何对象：剧烈分化日的点远离原点，普涨普跌日的点沿对角线聚集。\par
\textbf{行业相关性点云}。一点一申万二级行业，坐标为该行业最近 10 期简单收益（由 11 期收盘价构造）；行业 i、j 间距离取 \(d(i,j) = \sqrt{2(1 - \rho_{ij})}\)，其中 \(\rho_{ij}\) 为两行业收益的相关系数。该距离把相关结构转为度量空间：行业高度共动时点云收缩，分化时点云展开。节点数随历史在 67—131 之间浮动（已披露）。\par
两类点云的拓扑摘要均取至 H2：市场状态点云取 H0/H1 的总持久性与相邻日威尔斯特拉斯距离；行业相关性点云取 H0/H1/H2 的持久性统计、景观范数，并单独记录二维空腔计数 N（券商规则核心量）。两类点云中的同名摘要不可交叉解释。实现基于开源拓扑工具箱 giotto-tda（Tauzin 等，2021；文献 [35]），距离与滤过参数在预登记后固定。\par
\section{第二章　量化投资领域拓扑数据分析文献评述}
本章先回答一个前置问题：拓扑方法在量化投资实践中目前走到了哪一步？一句话总结：有应用价值的雏形已经出现，但经过严格验证、可落地复用的过硬证据依然不足。我们先梳理现有研究的整体版图，再分析这类方法的独特优势，同时明确其存在的缺陷与证据边界，最后给出对整个领域证据强度的整体判断。\par
\subsection{2.1　文献版图：五个研究方向}
我们共梳理了 94 篇拓扑方法在量化投资领域的应用研究（检索过程与逐条结论汇总见附录），按研究主题分为五个方向。为了方便对照，我们按统一口径给每个方向的证据强度分级：\textbf{A 级}指经过跨市场验证、由独立团队复现、计入真实交易成本后结果依然稳定；\textbf{B 级}指有规范的样本外检验，或方法上较严谨，但市场覆盖范围、可用事件数量或独立复现程度有限；\textbf{C 级}指仅来自单一研究、对比基线偏弱，或缺少严格的验证环节；\textbf{D 级}仅为概念性的初步探索。需要提前说明的是，整个领域目前没有达到 A 级的研究方向。\par
\begin{enumerate}
\item \textbf{危机预警与市场状态识别（B 级）}：这是目前文献数量最多的方向。从 Gidea 与 Katz 的开创性工作开始，研究者用持久景观的几何变化捕捉市场状态切换；Guo 等将这类方法扩展到多市场场景，发现零维特征（连通分支的变化）的预测能力可能比一维同调更强；Aguilar 与 Ensor 将拓扑指标与均值回归模型结合，Ismail 等把拓扑指标与临界慢化指标并列对比，同时报告了告警的精确率与误报率；Baitinger 与 Flegel 构建了拓扑湍流指数并完成了样本外检验；Guritanu 等则给出了带因果顺序的决策规则，是目前检测流程设计最规范的研究。称这个方向“最成熟”，指的是研究范式已经基本成型，并不代表其预测能力已经经过实盘检验、形成确定结论。
\item \textbf{波动率与尾部风险（B\(-\) 级）}：这类研究将持久同调与金融尾部特征的概念结合，把拓扑结构变化接入已有的波动率预测模型，在动荡期确实改善了预测效果（\textit{Journal of Finance and Data Science}, 2023），但其边际贡献是否仅为波动率信息的替代表达，还需要外部复现验证。
\item \textbf{组合构建与指数跟踪（B\(-\) 级）}：用拓扑特征辅助聚类、挑选代表资产、构建稀疏指数跟踪组合，一项覆盖 23 年样本的指数跟踪研究比较有说服力，但普遍存在对比基线不统一、交易成本处理口径不一致的问题。
\item \textbf{方向预测与策略交易（C\(-\) 级）}：将拓扑特征输入机器学习模型预测次日涨跌、直接构造交易策略的研究近年明显增多，但存在共同的方法缺陷：没有和非拓扑基线做严谨对比、不控制交易成本、调参过程不设隔离、不做多重检验校正——策略表现优异的原因未必来自拓扑方法本身。
\item \textbf{方法与工具层}：包括持久景观的统计化方法（Bubenik, 2015）、领域入门综述、以及开源软件工具箱（giotto-tda, 2021），这些基础工具为上述应用研究提供了底层支撑。
\end{enumerate}
\subsection{2.2　拓扑方法的独特优势}
相比传统统计方法，拓扑数据分析方法在金融场景的价值主要体现在五点：\par
\begin{enumerate}
\item \textbf{天然多尺度，无需人为设定阈值}：持久同调完整记录了不同尺度下“洞”的生灭全过程，并按特征存活时长自动加权，避免了传统网络分析中人为选择阈值带来的任意性。
\item \textbf{抗噪、抗形变}：第一章介绍的稳定性定理从数学上保证，输入数据发生连续小扰动时，拓扑摘要不会被放大成虚假的结构突变，非常适合含大量噪声的高频金融数据。
\item \textbf{能捕捉高阶结构}：环、空腔这类拓扑结构刻画的是相关矩阵无法直接呈现的多资产联动关系，比如轮动闭环、板块协同空腔，这些结构在线性统计里很难被直接观测。
\item \textbf{跨市场通用性强}：方法本身不依赖特定市场的微观交易结构，在美、欧、中、新加坡、印度、土耳其等不同市场都有成功应用的案例。
\item \textbf{对危机前后的结构变化敏感}：多数研究都观察到，在市场临近或处于危机的阶段，拓扑摘要会出现明确的结构性变化——这种对市场整体状态的刻画能力，是其作为市场状态识别工具的核心价值来源。
\end{enumerate}
\subsection{2.3　现有研究的缺陷与证据边界}
在肯定优势的同时，也必须明确现有文献普遍存在的问题：\par
\begin{enumerate}
\item \textbf{复现稳定性差}：同一个拓扑指标换一个市场、换一段样本，得到的结论可能完全相反；慢跌行情下信号往往明显滞后（Valdivia, 2023; Ruiz-Ortiz 等, 2022）。文献内部至少存在五组尚未达成共识的矛盾：\(H_0\) 和 \(H_1\) 哪个维度的信息更有效、范数该升还是降、反映的是内生危机还是外生冲击、拓扑信息是否只是波动率换了形式、不同的嵌入距离参数哪个更合理。
\item \textbf{参数自由度过高}：窗口长度、嵌入维数、滤过方式、同调维数、摘要统计量等可调参数构成了庞大的模型选择空间；多数研究不做消融实验与敏感性分析，若事后依据显著性选择参数，会产生较高的过拟合风险。
\item \textbf{仅在已知危机样本上展示信号，存在“检察官谬误”}：以事后已知事件筛选或展示信号会产生选择偏差，从而高估预警效果；必须严格区分“危机检测”和“危机预警”：危机期间指标出现显著变化，不代表该指标能在危机发生前发出信号，同期或滞后的响应不能视为提前预警。
\item \textbf{经济价值评估不足}：多数研究止步于统计指标层面的结果（精确率、提前量），很少把信号接入完整交易组合，在扣除误报导致的踏空成本与交易成本后进行严格的损益评估；即便是 Baitinger 与 Flegel 的样本外对比这样的工作也属少见，至于 Guritanu 等，则明确说明其检测器并未作为实盘交易策略做完整评估。
\item \textbf{可能只是波动率的复杂代理}：零模型预测研究显示，部分拓扑信号在对波动率做尺度标准化后，区分度会明显衰减（Rudkin, 2023）——我们自己的独立复现也观察到了类似现象。
\end{enumerate}
\subsection{2.4　证据强度总评}
把所有方向的证据强度做一个排序：市场状态/湍流识别为 B 级；危机早期预警为 B\(-\)；已实现波动率预测为 B\(-\)；稀疏指数跟踪为 B\(-\)；股票聚类/组合构建为 C+；方向预测/交易为 C\(-\)；衍生品直接定价为 D 级。整个领域目前没有 A 级证据——还没有出现跨市场、独立团队、扣成本后经过多期重复验证都成立的一致结论。\par
这正是本文实证工作的出发点：我们不复述“拓扑可以预测危机”的正面叙事，而是用事先登记的实验严格检验它的增量贡献，用完整的历史重放损益核算评估它接入组合后的边际经济价值。\par

\section{第三章　拓扑摘要用于危机预测的增量检验：未显示稳定增量}
从文献梳理可以看到，以 \(H_0\)、\(H_1\) 总持久性、相邻持久图间的威尔斯特拉斯距离为代表的拓扑摘要，在美股、加密货币市场的相关测试中普遍取得了较好的效果。自然可以提出假设：这类拓扑摘要应用到 A 股市场的危机预测场景，应该也能取得稳定的预测效果。为了验证这一猜测，同时避免做无价值的重复实验，我们专门设计了一组对照实验，核心要回答一个问题：在常规风险指标已经提供的预测能力之外，加入拓扑摘要能不能稳定提升危机预测效果？\par
\subsection{3.1　实验设计}
为避免事后调参与选择性报告，我们在接触数据之前即完整记录实验设计与判定标准（预登记），并明确停止规则：未达到预设判定标准即终止该研究路线，不根据结果反向修改参数以追求统计显著性。\par
实验的预测目标是未来 20 个交易日的二元危机标签（定义与全文研究口径完全统一）：当未来 20 个交易日内的最低累计收益跌到 \(-10\%\) 及以下时，标记为正例。\par
我们设置两个对照模型：\par
\begin{itemize}
\item \textbf{基线模型 M0}：仅使用五项常规风险指标，包括市场 50 日已实现波动、下行半离差、行业平均相关、行业协方差矩阵第一特征值（取最大特征值原值，非占比）、行业横截面离散度；模型采用固定 L2 惩罚的逻辑回归（正则化参数 \(C = 1.0\)），正负样本按权重平衡，固定随机种子 20261107。
\item \textbf{拓扑增强模型 M1}：在 M0 的基础上加入四个拓扑摘要特征（市场状态点云 \(H_0/H_1\) 的总持久性、相邻持久图间的威尔斯特拉斯距离等）。所有标准化过程仅在训练段拟合，用滚动前推（walk-forward）的方式计算样本外口径，完全杜绝未来信息泄漏。我们的核心判断指标是样本外精确率—召回率曲线下面积的增量（该面积衡量模型在控制误报的同时找出危机的能力，越高越好） \(\Delta =\) 面积(M1) \(-\) 面积(M0)，共做 17 折样本外测试。预先设定的“判定标准”非常严格：不仅要求 \(\Delta > 0\)，还要求 95\% 分块自助置信区间（块长 20、重抽 1000 次、种子 20261107）的下限也必须大于 0。
\end{itemize}
\subsection{3.2　结果：未达到预设判定标准}
16 个行业的全样本面板共 4070 个样本外交易日，包含 422 个正例日（正例率 10.4\%），样本外区间覆盖 2008-02-14 至 2024-12-03。核心结果如下表：
\begin{center}\small
\begin{tabularx}{\textwidth}{>{\raggedright\arraybackslash}X>{\raggedright\arraybackslash}X>{\raggedright\arraybackslash}X}
\toprule
\textbf{指标} & \textbf{M0（常规风险指标）} & \textbf{M1（加拓扑摘要）} \\
\midrule
精确率—召回率曲线下面积 & 0.1284 & 0.1209 \\
增量 \(\Delta\)（M1 \(-\) M0） & — & \textbf{\(-\)0.0076} \\
95\% 分块自助区间 & — & [\(-\)0.0207, +0.0143] \\
\bottomrule
\end{tabularx}
\end{center}
按行业拆分后结果同样不稳：13 个有效折里 4 个 \(\Delta\) 为正、8 个为负、1 个为零；27 个行业稳健面板的增量仅为 +0.0043，同样跨零；31 个行业的子样本里只剩 1 个正例，已经不具备统计解读价值。如果按实盘口径直接评估 M1 模型：事件召回率仅 0.125，年误报 5.4 次/年，告警精确率仅 0.032，中位提前期仅 1 天。\par
按照实验开始前预先固定的判定规则，这条路线未达到判定标准——我们按约定终止研究，不再继续增加参数或调整模型。\par
\subsection{3.3　如何解读这个负面结果}
首先要明确，置信区间跨零并不意味着本次检验证明了“拓扑摘要没有任何增量信息”——它只是说明在当前的实验设计下，我们的检验精度有限，小幅正增量和小幅负增量都在合理的可能范围内。因此本章的结论有非常严格的边界：我们只能说，在这套预先登记的设计下，直接拿拓扑摘要预测危机的路线没有显示出稳定的预测增量；不能把这个结论解读为“拓扑摘要完全不包含市场信息”。\par
这一结果和后续第五章“告警触发状态下拓扑特征有条件判别能力”的结论并不矛盾：合起来可以得到一个更准确的判断——拓扑摘要确实是带有状态信息的，但它更适合在市场已经进入告警区间后做精细的状态判别，而不适合作为独立的危机预测信号。\par
\textbf{本章小结}：拓扑摘要确实包含市场状态信息，但在预登记的严格实验设计下，作为独立的危机预测信号，它没有显示出稳定的显著优势。\par
\section{第四章　兴业证券研报策略的复现与检验}
\subsection{4.1　复现对象与口径}
原研报提出的可执行规则非常明确：基于申万二级行业的相关网络计算持续同调，提取二维空腔计数 N，当 \(N \le 1\)（即报告所称“市场孔洞少”）时进场做多，按原规则信号延续 4 个交易日，并称该规则可以获得可观的年化收益。这是中文实务文献里少见的、直接给出明确阈值与操作方向的 TDA 落地尝试，也是我们选择复现并检验它的原因。\par

这里需要先说明复现的边界：原文主结果使用的就是 Wind 渠道的申万二级行业分类，另以万得四级分类做稳健性检验（这部分四级分类的稳健性结果我们无法复现）；我们真正未能获取的是原文所用的 Wind 逐日行业名单与收益明细数据，因此本文用公开渠道的申万二级数据重建行业相关网络——行业层级与原文保持一致，差异仅在数据来源而非分类口径，不声称严格复现原文的回测数值，所有结论仅针对该规则的核心思想与可执行性，不对原文的历史回测结果做全盘否定。\par

复现按实盘交易时序设定：信号在 t 日收盘后计算，t+1 日开盘入场；信号日信息系数诊断的主评价期限为 t+1 开盘起算 20 个交易日的累计收益。整个复现的样本区间为 2015-07-01 至 2025-05-30。\par

\subsection{4.2　日频信息系数：点估计为正}
我们先计算信号与未来收益的日度信息系数，即二者之间的相关程度：对连续信号 \(s = -N\)（孔洞越少分值越高）和二值信号 \(1\{N \le 1\}\)，分别计算不同持有期下的信息系数（因为没有资产横截面，这里是纯时序序列信息系数，配对移动分块自助检验，块长 20、重复抽样 1000 次）。核心结果如下：
\begin{center}\small
\begin{tabularx}{\textwidth}{>{\raggedright\arraybackslash}X>{\raggedright\arraybackslash}X>{\raggedright\arraybackslash}X>{\raggedright\arraybackslash}X}
\toprule
\textbf{持有期} & \textbf{连续信号（\(-N\)）信息系数} & \textbf{二值信号（\(N \le 1\)）信息系数} & \textbf{备注} \\
\midrule
1 日 & 0.019 & — & 区间跨零 \\
5 日 & 0.029 & — & 区间跨零 \\
10 日 & 0.029 & — & 区间跨零 \\
20 日 & 0.066 & \textbf{+0.091} & 二值信号共 389 个信号日，信号后 20 日平均收益 +1.61\%，相对无条件收益基准 +0.26\%，胜率 64.8\%；主评价期自助区间跨零 \\
40 日 & 0.091 & \textbf{+0.125} & 该格 95\% 自助区间 [0.010, 0.269]，不跨零 \\
\bottomrule
\end{tabularx}
\end{center}
在完整的 50 格“信号—期限”矩阵里，有 4 格的自助区间不跨零。这里需要做一个诚实的判断：日度平均信息系数的点估计稳定在 0.09—0.13 区间，方向上确实支持“低孔洞数对应未来收益较好”的直觉，我们不能简单否定这个规则的信息含量；但同时要明确，主评价期（20 日）的自助区间是跨零的，统计上并不显著——稀疏信号的 IC 估计精度有限（信号出现次数少，有效样本量不足），单看 IC 点估计不足以直接判定规则有效。\par

\subsection{4.3　可交易性检验：三项严格口径均未通过}
“拿到规则就能直接交易”属于对这类研报策略的可交易性推断。本文从以下三个更严格的口径检验该推断：\par
\begin{enumerate}
\item \textbf{信号首次触发口径下信息系数不再显著}：实际交易中我们并不是每个交易日都做一次开平仓决策——一波趋势行情里信号会连续触发，本质是同一个交易机会被重复计数。如果仅保留每段连续信号的首次触发再重新计算，信息系数会降至 0.013，置信区间跨零，说明日度 IC 受到簇内重复信号的影响。
\item \textbf{固定阈值缺乏跨期稳定性}：原报告给出的固定触发阈值 \(N \le 1\)，在 2020 年 5 月之后就几乎不再触发，2021—2025 年整个区间内总共只有 13 个信号日。这个阈值并不是跨周期稳定的规律，其表现依赖特定样本阶段的市场结构——把 2015—2020 年的阈值直接套用到后续市场环境，规则便不再有效。
\item \textbf{相对常规超跌信号无增量预测能力}：“低孔洞数是好的入场点”这一说法，本质上和传统的超跌反弹信号高度重合。我们把它和普通的跌幅分位类超跌信号做对照后发现，加入孔洞数信息并没有带来额外的预测增量——超跌本身已经可以解释大部分收益，孔洞数并没有提供独立于传统技术信号的新信息。
\end{enumerate}
我们还做了分布稳定性检验：随着覆盖行业数扩容，N 的分布与节点数量持续漂移（行业数从 67 扩到 131），不同样本阶段中“N 数值低”并非同一统计事件；前期检验中“高 N 预示危机”这一反向假设也已经被本文实测证伪，为避免重复检验已被否定的假设，本轮不再对 N 的方向做新的验证。\par

\subsection{4.4　本章结论}
该规则包含一定的预测信息，日度频率下的表现不能完全否定，但作为独立的实盘交易信号，它缺乏跨期稳定性，不具备直接可交易性。这里需要把两个概念严格拆开：未获支持的是“按信号首次触发计算仍有效、固定阈值跨期稳定、相对超跌信号具有增量预测能力”这三项可交易性推断，而不是说规则本身完全没有信息含量。\par

这种分层评估框架也是本文对所有拓扑类规则的统一判读标准：不因信号具有统计信息就直接推断其可交易，也不因其未通过可交易性检验就否定其全部信息含量，而是明确区分“信号携带信息”和“信号能够作为独立交易策略获取收益”两个不同命题。\par
\section{第五章　危机状态判别模块的两方案设计与实证检验}
第三章的预登记实验未支持“单靠拓扑摘要预测危机发生时点”的路线，第四章对公开研报规则的复现也表明，单一拓扑阈值不能直接构成稳定的可交易策略。本章因此转向另一条路线：\textbf{不预测危机哪一天开始，而是在告警出现后逐日检查市场状态，并在证据充分时切换风控预案}。\par
\textbf{设计路径：先把金融规律翻译成可检验的指标，再固化规则。}本章并非从大量拓扑指标中事后挑出表现最好的一组。研究首先梳理金融文献中的共性结论（文献见附录）：熊市中资产相关性往往上升、分散化效果下降；第一主成分解释方差占比上升可能反映系统性风险积聚；危机还常伴随波动聚集、行业轮动减弱与风险沿网络扩散。随后将这些结论翻译成拓扑与相关结构指标，再通过预登记实验检验。事后分型得到方向相反的两组（15／30），其后的类型交互检验有四项通过假发现率控制（见 5.1）；与此同时，“危机中所有拓扑量同步跳升”的朴素预期未获支持。经过检验的方向最终被工程化为 5.2 节的七项状态确认指标与三项类型指标；触发层的候选信号同样来自短期急跌、异常波动和显著回撤等经典危机标志，再按预登记的识别与误报准则筛选（见 5.3）。因此，本章遵循的是“金融规律—指标翻译—事前检验—工程实现”的顺序；这些金融含义仍属事后机制解释，不构成因果验证。\par
\textbf{先用直白的话说明系统每天做什么。}收盘后，系统先检查短期急跌、异常波动和明显回撤三项条件；至少两项成立，才进入告警状态。告警不等于立即交易，系统还要继续检查市场相关结构与拓扑形态：只有连续两天得到一致的类型判断，且异常程度达到门槛，才切换仓位、止损或对冲方案。连续五天没有新信号便撤销告警，随后冷却十个交易日，避免信号反复出现时重复调仓。\par
整套模块自上游到下游共五层：
\begin{center}\small
\textbf{三项条件取二的触发规则（S3）}：是否开始警惕 \(\rightarrow\) \textbf{五日撤销、十日冷却的告警管理（B10）}\\
\(\rightarrow\) \textbf{危机状态确认}：异常程度是否足够 \(\rightarrow\) \textbf{危机类型投票}：当前形态更接近哪一套预案所对应的类型\\
\(\rightarrow\) \textbf{预案切换与历史重放损益核算}：采取什么动作，以及事后如何对账
\end{center}
S3、B10、C1、C2、R 均为研究档案中的版本编号。正文保留这些编号，是为了与预登记记录一一对应；阅读时只需先记住编号前的功能名称。\par
本章给出两套并行方案。\textbf{方案一是预登记后冻结的基准模型}：三项条件取二的触发规则＋五日撤销、十日冷却的告警管理＋7 项状态确认指标中至少 4 项达到极端条件＋3 项类型投票＋深慢／浅快两类风控预案，以及历史重放损益核算。\textbf{方案二是确认层扩展候选（研究档案记为 C1）}：只把状态确认扩展为 10 项指标中至少 6 项达到极端条件，其余模块与方案一完全一致。全文摘要、第六章与结论均以方案一为主口径；方案二作为重点候选同步呈现，不宣称已通过新的独立样本外验证，也不宣称可替代冻结基准。本章另记录一项只改变观察期起算方式的候选修订——\textbf{每段告警重新开放 20 个交易日观察期}（研究档案记为 R，原称“滚动判定钟”）：原规则只从事件归属时间范围（即允许告警计入该事件的日期区间，以下简称归属窗）内首个告警日起算一次 20 日观察期，R 改为每段连续告警各起算一次，其余组件不变。R 与两套确认层的组合均以候选身份呈现，其样本内择优性质在 5.5 节披露，不改变方案一的冻结基准地位。其余早期版本、无状态管理版本及未采纳候选只用于选型与消融对照，以保证研究过程可追溯。\par
各节安排如下：5.1 说明危机类型标签从何而来、证据有多强；5.2 说明告警之后每天如何判断、何时切换预案；5.3 说明何时开始告警、何时撤销与冷却；5.4 报告历史重放损益及拓扑特征的条件增量；5.5 集中列明现有证据不能支持的结论。\par
\subsection{5.1　危机状态事后分型与独立交互检验}
本节回答三个问题：危机类型标签如何构造、两套独立分型是否相互印证、拓扑指标的响应方向是否真的随危机类型而变。\par
\textbf{（一）类型标签的构造。}模型将危机状态划分为两类（为本文定义的研究标签，并非行业通用危机分类）：以市场状态点云整体尺度放大为主导的一簇（记为深慢型）与以行业相关结构松弛为主导的一簇（记为浅快型）。\textbf{该命名仅为标签约定，不表示两簇在财务上构成“跌得深、跌得久”与“跌得浅、跌得快”的二分}；划分依据是两簇在危机期间的主导拓扑变化方向，最大回撤与持续时长只作为事后画像。标签来自 45 个已标记危机事件的 11 项拓扑变化量：每项均取“危机段均值减去起点前 20 日均值”，其中市场状态点云类指标 6 项、行业相关网络类指标 5 项，两类点云的同名指标不做合并；所有变化量在 45 个事件上做标准化处理，聚类过程中的缺失值以该列中位数填补，最终采用欧氏距离与 Ward 层次聚类（逐步合并最相近的事件组）生成标签。\par
聚类数不是事后按金融结果挑选，而是在预登记的 \(k=2\) 至 8 中按轮廓系数（衡量组内相似、组间分离程度的指标）选择：\(k=2\) 的轮廓系数为 0.326，略高于 \(k=3\) 的 0.323，因此主分型保留两组，得到 15／30 两类。需要明确的是，0.326 只表示中等偏弱的聚类分离度，不代表市场天然存在边界清晰的危机二分结构；其他 \(k\) 的分组结果仅作稳健性对照，不进入模型标签体系。两组的事后金融画像如下：\par
\begin{center}\small
\begin{tabularx}{\textwidth}{>{\raggedright\arraybackslash}X>{\raggedright\arraybackslash}X>{\raggedright\arraybackslash}X>{\raggedright\arraybackslash}X>{\raggedright\arraybackslash}X}
\toprule
\textbf{拓扑主导方向（产品命名）} & \textbf{事件数} & \textbf{最大回撤中位数} & \textbf{持续日数中位数} & \textbf{事后 20 日收益中位数} \\
\midrule
深慢型（尺度放大主导） & 15 & \(-\)7.57\% & 15 & \(-\)1.56\% \\
浅快型（相关松弛主导） & 30 & \(-\)3.97\% & 13.5 & \(-\)6.16\% \\
\bottomrule
\end{tabularx}
\end{center}
上表为两簇的事后金融画像，两组数字为描述性中位数，未作组间显著性检验；分型依据不是这些金融量，回撤深度与持续日数的差异不应读作类型定义。\par
\textbf{（二）纯金融特征的独立复核。}为独立检验“危机类型是否调节拓扑指标的响应方向”，我们仅使用危机段最大回撤、持续交易日、事后 20 日收益三项纯金融特征重新做层次聚类，完全不引入任何拓扑输入。基于纯金融特征的独立分型得到 9 个深慢型、36 个浅快型事件，与拓扑聚类结果在 33/45 个事件上一致，分型一致率为 73.3\%。\textbf{这一基于纯金融特征的分型仅用于机制交互检验，不作为模型运行时的额外输入，也不另建第三套产品标签。}\par
两套事后分型在 45 个事件上的对应关系如下（括号内为占该拓扑类型的比例）：\par
\begin{center}\small
\begin{tabularx}{\textwidth}{>{\raggedright\arraybackslash}X>{\centering\arraybackslash}X>{\centering\arraybackslash}X>{\centering\arraybackslash}X}
\toprule
 & \textbf{金融深慢（9）} & \textbf{金融浅快（36）} & \textbf{合计} \\
\midrule
尺度放大主导（拓扑，15） & 6（40\%） & 9（60\%） & 15 \\
相关松弛主导（拓扑，30） & 3（10\%） & 27（90\%） & 30 \\
\midrule
合计 & 9 & 36 & 45 \\
\bottomrule
\end{tabularx}
\end{center}
两套分型在恢复路径上的差异同样可以核算。以万得全A 收盘价计，以事件窗内最低收盘价为起点、事件起点收盘价为目标，在随后 250 个交易日内首次收盘回到目标价视为恢复：金融深慢型 9 个事件中 5 个恢复（56\%），已恢复者中位 67 个交易日；金融浅快型 36 个事件中 27 个恢复（75\%），中位 52 个交易日。窗口结束后 20 个交易日，深慢型累计收益中位数为 +3.75\%，浅快型为 \(-\)6.27\%，即前者已出现反弹、后者仍在延续下跌，但深慢型的完全恢复比例与恢复速度均低于浅快型。\par
\begin{center}\small
\begin{tabularx}{\textwidth}{>{\raggedright\arraybackslash}X>{\centering\arraybackslash}X>{\centering\arraybackslash}X}
\toprule
\textbf{口径} & \textbf{金融深慢（9）} & \textbf{金融浅快（36）} \\
\midrule
事件后 20 日累计收益（中位数） & +3.75\% & \(-\)6.27\% \\
250 个交易日内回到起点价（恢复事件数 / 占比） & 5 / 56\% & 27 / 75\% \\
已恢复者所需交易日（中位数） & 67 & 52 \\
\bottomrule
\end{tabularx}
\end{center}
这里必须严格区分“事后评价标签”与“实时判别结果”：上述 15/30 的拓扑聚类标签与 9/36 的金融聚类标签都是事后复盘用于分组分析的参照，\textbf{73.3\% 衡量的只是这两套事后分型之间的对应关系，与实时判别无关}。模型每日运行时无法预知事件最终的类型归属，仅通过 5.2 节的三项指标做实时分类；实时判型与事后评价标签的匹配情况另在 5.4 节的已切换事件子集上报告，不能解释为全样本实时分类准确率或独立样本外分类准确率。\par
\textbf{（三）类型交互检验。}我们要回答的是：同一项拓扑指标在深慢与浅快两类事件中，是否呈现不同的变化方向。为此，在 9/36 的金融分型上对八项拓扑变化量做两组置换检验。这里的“置换”就是反复打乱 45 个事件的类型标签，再看真实分组产生的差异是否明显大于随机分组。检验单位是事件而不是交易日，统计量为“深慢组中位数减浅快组中位数”，共置换 10000 次计算双侧 \(p\) 值；置信区间另按事件整体自助抽样 5000 次。由于同时检验八项指标，统一采用 BH 假发现率控制（\(q<0.10\)），降低个别指标偶然显著造成的误判。最终四项通过校正，且两组的变化方向均相反：前三项为深慢型上升、浅快型下降；第四项（相关网络零维死亡均值）则为深慢型下降、浅快型上升。\par
\begin{center}\small
\begin{tabularx}{\textwidth}{>{\raggedright\arraybackslash}X>{\raggedright\arraybackslash}X>{\raggedright\arraybackslash}X>{\raggedright\arraybackslash}X>{\raggedright\arraybackslash}X}
\toprule
\textbf{拓扑变化量} & \textbf{深慢型中位数} & \textbf{浅快型中位数} & \textbf{置换 \(p\) 值} & \textbf{组间差异效应量（克利夫）} \\
\midrule
市场状态窗口均方根 & +0.0166 & \(-\)0.00175 & 0.0382 & +0.333 \\
市场状态平均两两相关 & +0.0769 & \(-\)0.02975 & 0.0219 & +0.568 \\
市场状态第一特征值占比 & +0.0705 & \(-\)0.0245 & 0.0291 & +0.543 \\
相关网络零维死亡均值 & \(-\)0.0735 & +0.02155 & 0.0127 & \(-\)0.444 \\
\bottomrule
\end{tabularx}
\end{center}
其余三项在两类中的中位变化同向；相关网络一维环数量虽呈反向差异，但未达到统计显著且检验功效有限。在更严格的 Bonferroni 多重检验校正下，0/8 通过；预登记已声明主口径采用上述假发现率控制。\textbf{因此，样本内证据支持的是“四项特定拓扑响应存在危机类型交互”，不能概括为所有拓扑指标都遵循同一方向，更不能外推为样本外预测能力。}\par
\subsection{5.2　日频危机状态确认与类型判别规则}
本节回答：告警之后，系统每天检查什么，满足什么条件才切换预案。两套方案都在触发后\textbf{逐日更新}两项结果：异常程度是否足够，以及当前形态更接近哪种类型；并不是等满 20 个交易日后才做一次分类。早期单特征诊断曾以 \(k=5\) 为主要检验窗口，那只用于候选指标筛选，不是最终产品的决策规则。三项类型指标的方向来自全样本交互检验，属于样本内知识的固化，并非只用 2015 年以前的数据独立估计得到。\par
\textbf{（一）方案一的危机状态确认评分（7 项指标投票）。}这一步只回答“当前是否足够异常”。对第 \(j\) 项指标定义一个 0/1 标记 \(I_{j,t}\in\{0,1\}\)：当天达到该项的极端条件记 1，否则记 0。确认得分与整数确认票分别为
\[
C_t^{(m)}=\frac{1}{m}\sum_{j=1}^{m}I_{j,t},\qquad K_t^{(m)}=\sum_{j=1}^{m}I_{j,t},\qquad m\in\{7,10\}
\]
计票规则如下：
\begin{enumerate}
\item 每项指标只与自身在\textbf{当日之前的 756 个交易日}中的历史水平比较，不把当日值放进自己的比较基准；
\item 正向指标高于历史 85 分位，记为达到极端条件；负向指标低于历史 15 分位，记为达到极端条件；
\item 历史有效观测不足 120 个，或当天指标缺失，该项直接记 0，不凭缺失数据投票；
\item 分母固定为 7 或 10，不改成“当天有数据的指标数”。这样规则始终可以手工复核，也避免数据缺失时分母变小、门槛被被动降低。
\end{enumerate}
方案一的七项指标按风险维度分组：\par
\begin{center}\small
\begin{tabularx}{\textwidth}{>{\raggedright\arraybackslash}X>{\raggedright\arraybackslash}X>{\raggedright\arraybackslash}X}
\toprule
\textbf{风险维度} & \textbf{指标及方向} & \textbf{风险含义} \\
\midrule
相关性上升（3 项） & 平均两两相关上升、首特征值占比上升、行业平均相关上升 & 资产同涨同跌，分散化效果下降 \\
拓扑结构简化（3 项） & 二维空腔计数下降、一维环计数下降、二维持久景观下降 & 行业轮动结构减弱，可分散风险的空间收窄 \\
风险传染（1 项） & 零维同调死亡尺度下降 & 连通分支更早合并，风险扩散门槛下降 \\
\bottomrule
\end{tabularx}
\end{center}
危机状态确认层\textbf{不直接包含波动率变量}，其信息来自相关性上升与拓扑结构变化；最终 S3 触发器包含波动率条件，二者为独立层级，不能混为一层。指标的金融含义均为事后机制解释，相关研究边界统一在 5.5 节列明。\par
\textbf{（二）方案二的确认层（10 项指标）。}方案二把确认标准由“7 项中至少 4 项极端”改为“10 项中至少 6 项极端”。新增的三项均来自市场状态点云：相邻日持续图距离上升、零维死亡值均值上升、一维环数量下降；十项指标沿用同一套历史分位数规则。需要明确的是，方案二同时改变了指标数量和确认门槛，因此不能把结果提升全部解释为新增拓扑特征的独立贡献。\par
\textbf{（三）危机类型判别（3 项指标投票）。}这一步回答“当前形态更接近哪一套预案对应的类型”，比较的不是某一天的高低，而是触发前后的变化：以触发日前 10 个交易日的均值为基准，与触发日至当前判定日的均值相比较。基准至少要有 8 个非缺失观测，当前窗口至少要有 80\% 的观测有效。市场状态窗口均方根上升、平均两两相关上升、相关网络零维死亡均值下降时，各投一票深慢型；反向则投浅快型。类型分达到 2/3 判为深慢型，不超过 1/3 判为浅快型，其余记为无法判定；有效类型指标不足 2 项时记为数据不足，不凭单项指标强行选择预案。\textbf{三项指标识别的是点云尺度、行业共动与相关网络的形态变化，不是预先判断事件最终会跌多深、持续多久。}\par
\textbf{（四）预案切换规则。}触发后的第 1—20 个交易日为类型判断的观察期。状态确认回答“异常是否足够”，类型投票回答“应选择哪套预案”；\textbf{两类条件必须同时满足才实际调仓}。因此，系统可以在异常不足时继续观察，也可以在类型分歧或数据不足时维持原仓位，而不是一旦出现告警就立即交易。具体条件是：同一类型连续 2 日得到确认，且第二日达到相应方案的状态确认门槛（方案一至少 4/7 项，方案二至少 6/10 项），方可在确认日收盘切换，新仓位自次一交易日起计损益；最多允许 1 次类型修正。历史重放还要求确认日早于事后已知的事件结束日，这一前视限制在 5.5 节单独披露。原规则只从首个告警日起算一次 20 日观察期，可能错过晚到的证据；相应瓶颈及候选修订见本节（六）。\par
历史重放中的切换统计两方案分别如下：\textbf{方案一}实际切换 19/45 个事件，首次切换相对触发日的中位延迟为 2 个交易日，第 25 至第 75 百分位为 2—4 日，完整范围为 2—20 日；\textbf{方案二}实际切换 17/45 个事件，中位延迟为 3 日，第 25 至第 75 百分位为 2—5 日，完整范围为 2—20 日。两个方案在 45 个事件的事件侧核算中均发生 3 次类型修正，说明纠错机制确实被使用，但并未引发频繁的反复切换。\par
\textbf{（五）两方案共用的预案参数。}这里的“深慢／浅快”只决定风控强度与工具组合，不产生方向性做多信号；无法判定或数据不足时维持原仓位，不将不确定状态强行映射为交易动作。\par
\begin{center}\small
\begin{tabularx}{\textwidth}{>{\raggedright\arraybackslash}X>{\raggedright\arraybackslash}X>{\raggedright\arraybackslash}X>{\raggedright\arraybackslash}X}
\toprule
\textbf{实时判型} & \textbf{目标仓位与止损} & \textbf{对冲安排} & \textbf{设计目的} \\
\midrule
深慢型 & 50\% 仓位；相对切换日收盘价回撤 20\% 止损 & 买入虚值认沽：执行价为入场价 0.95 倍，名义仓位 0.5，权利金率 1.5\% & 降低长期回撤暴露，同时保留尾部保护 \\
浅快型 & 80\% 仓位；相对切换日收盘价回撤 12\% 止损 & 不配置认沽 & 保留较高仓位，以较紧止损控制短期冲击 \\
无法判定／数据不足 & 维持原仓位与原止损 & 不新增对冲 & 等待更多证据，避免在信息不足时误切预案 \\
\bottomrule
\end{tabularx}
\end{center}
止损监测在切换后的持有路径上执行；事件结束后按 20 个交易日线性恢复至满仓。深慢型认沽的权利金在切换时计入成本，到期赔付按预登记规则进入单位净值路径；因此后文的 +79.0 或 +104.9 并非只比较仓位，还包含止损与对冲现金流。\par
\textbf{（六）候选修订：每段告警重新开放 20 日观察期（记作 R）。}原规则只在第一次告警后给出一次 20 个交易日的观察期。如果观察期结束时证据仍不充分，后面即使再次告警，也不会重新计时。本小节讨论能否只改这一处时序规则，同时保持其他冻结组件不变。\par
对“已经告警但最终没有切换”的事件逐日复算后，可以看到两类不同的原因。第一类是\textbf{证据始终不够}：整个危机期间，达到极端条件的指标数从未达到门槛，增加观察机会也无济于事，本文称之为“水平型”瓶颈。第二类是\textbf{证据到得太晚}：确认条件后来确实满足，但 20 日观察期已经结束，本文称之为“时钟型”瓶颈。后一类与 S3 偏向提前告警的设计有关：告警可能早于主跌段，而观察期又从归属窗内第一个告警日起算，主跌段来迟时，证据与观察期就会错开。\par
2022 年事件 12 是完整例子。归属窗内首个告警日为 2022-01-25，20 日观察期至 2022-03-01 结束时，确认指标仅为 0—1 项；真正达到门槛的 10 个交易日落在 2022-03-15 至 03-29，对应俄乌冲突与局部封控下的主跌段，确认指标为 4—5 项，全部位于观察期之外。2008 年 1 月（事件 117）与 2016 年 1 月熔断（事件 2）也有类似情况：第二段连续告警内本有满足切换条件的日期，但原规则不重新开放观察期，因而未切换。\par
R 只修改这一处：\textbf{每段与事件归属窗相交的连续告警，都各自重新起算一个 20 交易日观察期}。切换条件——同型连续 2 日、第二日确认指标达到门槛、确认日早于事件结束日、最多 1 次类型修正——以及 S3、B10、确认指标、类型投票和预案参数全部不变。候选实验按升版预登记执行，登记先于计算，并事先写明由诊断导出的预期变动清单（三个新切换与两个类型修正，见 5.4（六））。实现时先禁用新规则进行自检：六组配置共 270 个事件行必须与修订前逐字段一致，通过后才运行实验配置，实验产物再经双跑核对。\par
这项修订有三条可以直接检查的等价关系，使每组配置的变动恰好收敛于五个事件：
\begin{enumerate}
\item \textbf{误报段不变}：原规则对误报段本来就逐段判断、分别起算观察期，因此修订前后的误报激活段数、逐段损益与误报合计完全相同；
\item \textbf{只有一段告警的事件不变}：归属窗内只有一段连续告警时，新旧规则相同；六组配置中每组其余 40 个事件的判定路径均逐字段复现修订前结果；
\item \textbf{验证期重放不变}：重放逐行逐列与修订前一致，5.4（七）的结论对 R 同样成立。2026 年三段“证据始终不够”的窗口在 R 下依然未切换，类型判别短板并未因此修复。
\end{enumerate}
因此，R 挽回的是“证据晚到、观察期却已结束”的事件，不能解决“证据始终达不到门槛”的事件。其样本内损益贡献、与方案二的组合结果及证据边界，分别见 5.4（六）与 5.5 节。
\subsection{5.3　触发层与告警状态管理：从基线到全样本危机识别}
本节回答：判别流程从什么时候开始，告警如何维持与撤销，以及 45/45 的全样本识别到底覆盖了什么。\par
\textbf{（一）三项条件取二的触发规则（S3）。}两个方案共同使用同一触发规则。S3 是研究档案编号，其含义就是下面三项条件中\textbf{至少两项同时成立}：
\begin{enumerate}
\item 最近 5 个交易日简单收益之和不超过 \(-\)4\%；
\item 20 日已实现波动率高于自身滚动 756 日历史的 70 分位，且历史有效观测不少于 250 个；
\item 收盘价距最近 20 日最高收盘价的回撤达到 5\%。
\end{enumerate}
三项分别检查短期下跌、相对自身历史的异常波动和近期价格位置。要求三中取二，而不是任意一项就告警，是为了避免单日下跌或单纯高波动直接主导系统。这里的五日收益按简单收益率相加，与第一章拓扑点云使用的对数收益属于不同模块的既定口径。该规则不估计危机概率，也没有训练得到的回归系数，每一步都能手工检查。研发早期的初版同样采用三条件取二，但第三项是跌破 120 日半年线；最终规则将其替换为 20 日高点回撤，这一历史差异仅用于说明选型来源。\par
\textbf{（二）五日撤销、十日冷却的告警管理（B10）。}触发之后，系统按“\textbf{连续 5 日无原始命中即撤销，撤销后冷却 10 个交易日}”管理告警状态。B10 是这组参数在研究档案中的编号。具体而言：
\begin{enumerate}
\item 原始 S3 命中使非告警状态进入告警状态；
\item 告警期间若再次命中，连续无命中计数归零；
\item 只有连续无命中达到 5 日才撤销告警，且第 5 日不再计为告警日；
\item 撤销后的 10 个交易日为冷却期，即使原始条件再次命中也不立即重启告警，冷却结束后才重新等待 S3 命中。
\end{enumerate}
全样本共发生 117 次撤销。“不提前撤销”是指不采用“静默 3 日且收复 250 日线”的提前解除分支，并非永不撤销。状态机只管理告警的持续、撤销与再次激活，不改变上述三项原始触发条件。\textbf{B10 的目的不是把告警日数压到最低，而是避免相邻原始命中和短暂重启造成重复预案激活；因此告警日占比上升与误报激活数下降可以同时发生。}连续两日同向确认是 5.2 节的预案切换条件，不属于状态机，也不存在“信号消失后保留十日缓冲”的规则。\par
\textbf{（三）全样本识别与预案启动口径。}这里要分开回答两个问题。第一，系统有没有识别到危机：统计时允许信号出现在\([\text{起点}-20\ \text{日},\ \text{结束}+20\ \text{日}]\)，也就是危机结束后仍留有 20 日补充观察期。第二，系统来不来得及实际启动预案：用于损益核算的切换资格只到事件结束日为止。在这一共同的触发与告警层下，识别口径达到 \textbf{45/45}：\textbf{24 个事件在起点前触发、7 个在事件窗内触发、14 个直到事件结束后的补充观察期才触发}。前两类合计 \textbf{31/45}，来得及启动预案核算；其余 14 个虽然仍计为已识别，但因事件已经结束，不进入危机事件内的实时预案损益核算。在 31 个可核算事件中，方案一切换 19 个、未切换 12 个，方案二切换 17 个、未切换 14 个。\par
两方案的告警日均占全历史 23.5\%，误报段均为 83 段。由于共用 S3 与 B10，上述识别、时点分档、告警日和误报段数字完全相同；差异只发生在危机状态确认、预案切换及其后的损益核算（方案一实际激活 33 个误报段，方案二为 21 个）。因此，阅读时应按“\textbf{45/45 识别 \(\rightarrow\) 31/45 可核算 \(\rightarrow\) 19/45 或 17/45 实际切换}”的顺序理解，而不能把全部识别解释为全部提前预警或全部实施了保护；告警日占比也不能直接解释为误报率或误报成本。\par
\textbf{（四）历史选型与消融（非产品口径）。}以下历史结果只用于说明选型与消融，不是另一套产品口径。早期原实验按归属窗内可核算口径，事件＋价格确认、区间检测与混合规则分别为 1/45、21/45 和 29/45；这些数字不是最终产品的宽识别窗召回率。随后候选信号扩展到 15 项，并围绕初版三条件取二规则进行并联、替换单项条件与表决方式的比较：初版规则在宽识别窗下为 44/45、告警日占比 21.1\%；原始 S3 在不接状态机时达到 45/45，分档为 26／6／13，告警日占比降至 17.3\%，识别窗外误报为 9.9 日/年。\textbf{17.3\% 因此只是“S3、无状态管理”的对照数字，不能与最终 S3＋B10 的损益结果拼接。}\par
状态机选型分两步。研发早期先在一版触发规则（研发代号 P8，早于最终 S3 定稿）上评估九组状态机变体：无治理时误报为 62.0 日/年；只加撤销而不冷却时，两类撤销规则分别升至 107.1 和 121.2 日/年，说明撤销可能合并并延长相邻告警，而不是自动减少误报；在相同撤销规则下增加冷却期，误报天数才逐步下降，但也伴随部分事件响应推迟。撤销规则本身只改变告警的持续与分段，不能稳定降低误报总量；上述数字属于该早期版本，不能混入最终 S3＋B10 的产品结果。此后在最终 S3 触发规则上重新比较无治理对照与八组状态机，按预登记的“45/45 识别硬门槛、净损益差、误报日、冷却长度”顺序选定 B10。\par
\subsection{5.4　历史重放下的预案损益与拓扑特征的条件增量}
本节回答两个问题：这套规则在历史重放中值多少损益；拓扑特征相对既有确认规则有没有产品侧增量。为避免把不同证据等级的结果混读，全文按三个层次组织：第一层是\textbf{冻结基准方案一}，回答事先冻结的规则在历史上做到了什么；第二层是\textbf{两个独立改进方向}，方案二检验增加确认指标并提高门槛能否减少误报后的实际调仓，R 检验每段告警重新给予 20 日观察期能否挽回“证据到得太晚”的事件；第三层是\textbf{样本内最优组合方案二＋R}，回答两个改进合并后历史上最高达到多少。第三层只是后续前瞻验证的首选候选，不改变第一层的全文主口径。\par
\textbf{（一）核算方法：基于历史重放的预案损益对照核算。}固定已经发生的历史价格，只改变假设采取的风控动作，从而比较“采取预案”与“不采取预案”的差异。每个事件或误报段都假设起始资金为 1.00，逐日复利，并计入止损、认沽对冲现金流和结束后 20 日线性回补。四条路径平行对照：\par
\begin{center}\small
\begin{tabularx}{\textwidth}{>{\raggedright\arraybackslash}X>{\raggedright\arraybackslash}X>{\raggedright\arraybackslash}X}
\toprule
\textbf{路径} & \textbf{设定} & \textbf{角色} \\
\midrule
事后正确预案 & 触发后第 2 日使用事后正确类型 & 机制上限 \\
反向类型预案 & 同一时点故意使用相反类型 & 反向对照 \\
无预案基准 & 不采取预案 & 基准路径 \\
实时判定切换 & 信号输入只用当时可得的确认票与类型票 & 实盘模拟结果 \\
\bottomrule
\end{tabularx}
\end{center}
实时路径的切换资格仍受事后已知的事件结束日约束，前视偏差在 5.5 节单独披露。上限与实时路径的差距反映确认延迟、未切换与判型错误造成的损耗；反向对照用于检查风控损益是否真的依赖类型方向。事件路径在事件结束后 20 日线性回补；误报段则统一采用“20 日持有观察＋20 日线性回补”的比较窗口。净损益差合计定义为
\[
\text{净损益差合计}=\sum_{e=1}^{45}\left(V_e^{\text{实时}}-V_e^{\text{无预案}}\right)
+\sum_{j\in\text{误报段}}\left(V_j^{\text{实时}}-V_j^{\text{无预案}}\right)
\]
未触发、未切换，或直到事件结束后的补充观察期才首次识别的事件，按构造贡献 0。每项结果以单位净值百分点报告，再在事件与误报段之间求算术和；\textbf{它用于比较预案与不采取预案的差异，不是同一账户沿时间累积形成的收益率，也没有处理重叠事件之间的资金占用。}\par
\textbf{（二）主结果表。}下表前两行分别给出冻结基准方案与确认层扩展候选（C1）；第三、四行为候选修订 R（每段告警重新开放 20 日观察期，规则见 5.2（六），组合评估见 5.4（六））与两套确认层的组合，属候选记录，证据等级低于前两行；其余行均为触发器／状态机消融或未采纳的候选增强实验。\par
\begin{center}\small
\begin{tabularx}{\textwidth}{>{\raggedright\arraybackslash}X>{\raggedright\arraybackslash}X>{\raggedright\arraybackslash}X>{\raggedright\arraybackslash}X>{\raggedright\arraybackslash}X}
\toprule
\textbf{配置方案} & \textbf{宽窗识别；可核算} & \textbf{告警日占比} & \textbf{净损益差合计（单位净值百分点）} & \textbf{说明} \\
\midrule
\textbf{方案一（冻结基准）：S3＋B10＋7项4票} & \textbf{45/45；31/45} & \textbf{23.5\%} & \textbf{+79.0} & 19 个事件切换；事件侧损益差 +147.4，误报段 \(-\)68.4 \\
\textbf{方案二（C1重点候选）：S3＋B10＋10项6票} & \textbf{45/45；31/45} & \textbf{23.5\%} & \textbf{+104.9} & 17 个事件切换；事件侧损益差 +153.1，误报段 \(-\)48.3 \\
候选修订组合：S3＋B10＋7项4票＋分段重新计时 R & 45/45；31/45 & 23.5\% & +114.8 & 22 个事件切换；事件侧损益差 +183.2，误报段 \(-\)68.4（与方案一逐分相同） \\
候选修订组合：S3＋B10＋10项6票＋分段重新计时 R & 45/45；31/45 & 23.5\% & +140.7 & 20 个事件切换；事件侧损益差 +188.9，误报段 \(-\)48.3（与方案二逐分相同）；证据边界见 5.5 \\
消融：旧混合触发规则 & —；29/45 & — & +27.0 & 原实验只报告事件结束前可核算数，不与上述识别口径直接比较 \\
消融：S3，无告警状态管理 & 45/45；32/45 & 17.3\% & \(-\)5.3 & 事件侧 +114.1，误报段 \(-\)119.3 \\
候选增强（未采纳）C2：拓扑指标恢复常态后提前解除告警 & 45/45；29/45 & 20.4\% & +30.2 & 误报日减少，但误报段及成本上升，净损益差低于对照配置 \\
\bottomrule
\end{tabularx}
\end{center}
表中“宽窗识别”允许信号出现在事件起点前 20 日至结束后 20 日；“可核算”则要求信号不晚于事件结束，且最早只追溯到起点前 20 日。旧混合规则的原实验未按这一识别口径重算，故以“—”标示，不把它在可核算口径下的 29/45 与最终方案在宽识别口径下的 45/45 直接比较。表中各行及本章后续各表的净损益差合计均由精确值四舍五入得到，与事件侧、误报段两个分项之差可能存在 0.1 的尾差。\par
\textbf{（三）如何读这张表。}从原始 S3 到 S3＋B10，事件侧损益差合计增加约 33.3，误报成本改善约 50.9，合计改善约 84.3（分项与合计分别四舍五入，二者相加差 0.1）；两项均有贡献，后者更大。全样本“实时路径减无预案路径”的中位数为 0，因为未触发或未切换事件按构造记 0；方案一的 19 个已切换事件上，该差值中位数为 +1.157，13 个为正、6 个为负；方案二的 17 个已切换事件上，中位数为 +3.220，13 个为正、4 个为负。净损益差应与已切换子集结果及召回分档共同阅读，不能只看被大量零值拉平的全样本中位数。\par
\textbf{（四）方案二相对方案一的差异。}方案二净损益差提高 25.8，其中事件侧全程收益差由 +147.4 增至 +153.1，仅改善 5.7；误报段损益由 \(-\)68.4 改善至 \(-\)48.3，贡献 20.1。误报激活段由 33 降至 21，实际切换事件由 19 降至 17，首次切换延迟中位数由 2 日增至 3 日。\par
这里需要说明为什么两方案告警日占比同为 23.5\%、误报段同为 83 段，误报损益却能不同：告警日、误报段和误报日数都是 S3＋B10 状态机的产物，两方案共用这一触发与告警层，故三者完全相同；而误报段只有在通过确认投票、实际切换仓位时才产生损益，未激活的误报段不发生交易，按构造损益差为 0。方案二更严的确认层（10 项中 6 票）在交易动作发生之前就把 83 个误报段中的更多段过滤掉（激活段 33\(\rightarrow\)21），告警状态依旧存在，但更多误报没有转化为实际仓位动作。进一步的核对显示，平均每激活段成本并未下降——方案一约 \(-\)2.07/段，方案二约 \(-\)2.30/段——\textbf{改善全部来自激活段数减少，而非每段成本变低}；析因网格中其余配置组亦呈同一规律（10 项 4 票激活 43 段、成本 \(-\)88.9；7 项 6 票激活 14 段、成本 \(-\)44.2）。\par
\textbf{因此，C1 的主要样本内优势是减少误报激活及其成本，而不是显著提高危机期间收益；代价是确认更慢、切换事件更少。}扣除各自最大正贡献事件后，方案二净损益差仍为 +51.3，方案一同口径为 +26.6。该检查说明改善并非完全由单一事件贡献，但它只是收益集中度描述，不构成统计显著性证明。\par
\textbf{（五）C1 补全实验：扩容与门槛的析因网格。}为拆分“特征扩容”与“确认门槛加严”两类效应，补充实验在完全相同的 S3＋B10 产品框架上预登记并运行“7／10 项指标 \(\times\) 4／5／6 整数票”六组网格。该设计固定所有其他产品组件，组间差异只能来自确认层及其后续损益路径；六组的 45/45 识别、24／7／14 分档、31/45 可核算、告警日与误报段均通过不变量校验。\par
\begin{center}\small
\begin{tabularx}{\textwidth}{>{\raggedright\arraybackslash}X>{\raggedright\arraybackslash}X>{\raggedright\arraybackslash}X>{\raggedright\arraybackslash}X>{\raggedright\arraybackslash}X>{\raggedright\arraybackslash}X}
\toprule
\textbf{确认配置} & \textbf{切换事件数} & \textbf{事件侧损益差} & \textbf{误报激活段} & \textbf{误报段损益差} & \textbf{净损益差合计} \\
\midrule
7 项 4 票（方案一） & 19 & +147.4 & 33 & \(-\)68.4 & +79.0 \\
10 项 4 票 & 20 & +154.0 & 43 & \(-\)88.9 & +65.1 \\
7 项 5 票 & 17 & +147.6 & 24 & \(-\)52.9 & +94.8 \\
10 项 5 票 & 18 & +155.1 & 30 & \(-\)71.9 & +83.1 \\
7 项 6 票 & 13 & +117.7 & 14 & \(-\)44.2 & +73.5 \\
10 项 6 票（方案二） & 17 & +153.1 & 21 & \(-\)48.3 & \textbf{+104.9} \\
\bottomrule
\end{tabularx}
\end{center}
按预登记的分解方式，可以把方案二的提升拆成三部分：只增加指标、只提高门槛，以及两项改动组合后额外产生的交互差异。同 4 票仅扩容至 10 项，净损益差变化为 \(-\)13.9；原 7 项由 4 票加严至 6 票，变化为 \(-\)5.5；二者组成 10 项 6 票后相对原产品净提高 25.8。分解式如下（分解项保留两位小数以使恒等式闭合）：
\[
\Delta\text{净损益差}_{\mathrm{C1}}
=\underbrace{-13.93}_{\text{同4票扩容}}
+\underbrace{-5.52}_{\text{7项加严至6票}}
+\underbrace{45.30}_{\text{组合交互}}
=+25.85
\]
\textbf{这说明扩容并非在所有门槛下都有效，加严也并非单调改善；C1 的结果来自特征集与确认纪律的组合，不能拆解为新增拓扑特征的独立因果效应。}7 项 5 票在不增加特征时也达到 +94.8，进一步说明确认纪律本身重要；但它只是预登记的次要敏感性组，不能因结果较好就事后反选为产品。六组结果均予保留，不按最高净损益差自动改写产品。\par
\textbf{（六）候选修订 R 与确认层的组合评估。}本小节回答三个问题：每段告警重新起算观察期能带来多少变化；它与方案二组合后是否优于各自单独使用；为什么组合成绩最高也不能直接成为新的产品基准。组合评估在同一 S3＋B10 产品框架上运行 \(2\times 2\) 主网格：确认层取方案一（7 项 4 票）与方案二（10 项 6 票），观察期起算方式取原规则（只起算一次）与候选 R（每段告警重新起算），其余组件全部冻结。四格的识别、告警日与误报段数字完全相同，差异只能来自观察期的起算方式。\par
\begin{center}\small
\begin{tabularx}{\textwidth}{>{\raggedright\arraybackslash}X>{\raggedright\arraybackslash}X>{\raggedright\arraybackslash}X>{\raggedright\arraybackslash}X>{\raggedright\arraybackslash}X>{\raggedright\arraybackslash}X}
\toprule
\textbf{配置} & \textbf{切换事件} & \textbf{事件侧损益差} & \textbf{误报激活段} & \textbf{误报段损益差} & \textbf{净损益差合计} \\
\midrule
方案一（原观察期规则） & 19 & +147.4 & 33 & \(-\)68.4 & +79.0 \\
方案一＋R & 22 & +183.2 & 33 & \(-\)68.4 & +114.8 \\
方案二（原观察期规则） & 17 & +153.1 & 21 & \(-\)48.3 & +104.9 \\
方案二＋R & 20 & +188.9 & 21 & \(-\)48.3 & +140.7 \\
\bottomrule
\end{tabularx}
\end{center}
读表可得四个事实。\textbf{第一，R 的贡献独立于确认层门槛}：两个主要配置（方案一与方案二）的变动事件完全相同、逐事件增量逐分相同，净损益差同步提高 +35.8——观察期起算方式的修改不依赖确认门槛的高低；析因网格其余四组在 R 下亦全部为正（+14.3 至 +35.8），不存在某一门槛下失效的反例。\textbf{第二，改善全部来自事件侧，且高度集中}：误报激活段与误报成本逐分不变（表中两列上下对齐），事件侧合计由 +147.4 增至 +183.2；但 +35.8 中有 33.5 来自两场崩盘开局——2008 年 1 月（事件 117，+23.3）与 2016 年 1 月熔断（事件 2，+10.2），2022 年事件 12 仅 +1.1（切换时点接近阶段底部，其后第二波下跌令认沽对冲勉强回本）。五个变动事件（两个主要配置完全相同）的明细如下：\par
\begin{center}\small
\begin{tabularx}{\textwidth}{>{\raggedright\arraybackslash}X>{\raggedright\arraybackslash}X>{\raggedright\arraybackslash}X>{\raggedright\arraybackslash}X}
\toprule
\textbf{事件（类型）} & \textbf{修订前} & \textbf{修订后} & \textbf{变动机制} \\
\midrule
117（2008-01，深慢） & 未切换，贡献 0 & +23.3 & 第二段告警（起于 2008-01-17）于 01-23 确认深慢 \\
2（2016-01，深慢） & 未切换，贡献 0 & +10.2 & 第二段告警（起于 2016-01-05）于 01-07 确认深慢 \\
12（2022，深慢） & 未切换，贡献 0 & +1.1 & 第二段告警（起于 2022-03-08）于 03-15 确认深慢 \\
114（2006-07，深慢） & 切浅快，+1.2 & +2.4 & 首段告警沿用原切换；第二段告警于 07-21 确认深慢，触发类型修正，终判转对 \\
119（2008-07，浅快） & 切深慢，+46.8 & +46.8 & 第二段告警于 08-18 判浅快，类型标签修正转对；但修正日前 20\% 止损已使仓位退出，实际损益不变 \\
\bottomrule
\end{tabularx}
\end{center}
五个变动事件与升版预登记时写明的预期清单完全一致，无任何计划外变动——这正是“先登记后计算”提供的事前约束。\textbf{第三，切换覆盖与判型一致性同向改善，但不构成判型能力提高的证明}：切换事件 19→22／17→20，首次切换延迟中位数维持 2 日／3 日，首次判型与评价标签一致数 10→13／9→12，最终预案一致数 7→12／6→11；新增一致主要来自三个新切换事件本就判对深慢型与事件 114 的修正转对，分子分母同步扩大，不能读成判型准确率的提升。\textbf{第四，样本内最优组合的定位是优先前瞻验证，而非替换基准}：方案二＋R 的 +140.7 是本文已评估组合中的最高值，但 R 与 C1 的改善均在同一段已知历史上计量，且验证期已在此前的补全实验中被描述性使用，两个候选各自择优的偏差在组合上叠加，见 5.5 节第 4 条。因此，该组合是后续前瞻验证的首选候选，而不是已经验证完成的最终产品；冻结基准与全文主口径仍为方案一。\par
\textbf{（七）验证期历史重放与判型短板。}这里的“重放”是指把冻结规则重新应用于已经发生的数据，只作描述，不构成新的独立样本外检验。两方案的触发与告警层均保持 6/6 事件识别、窗外误报约 3.8 日/年及 2026 年零误报；但方案二未显示判型修复。样本内已切换事件上，首次判型与产品评价标签的一致数为方案一 10/19、方案二 9/17；最终预案一致数分别为 7/19 和 6/17，不能据此宣称判型准确率提高。三段叙事事件的重放细节如下，两个方案结果相同：\par
\begin{center}\small
\begin{tabularx}{\textwidth}{>{\raggedright\arraybackslash}X>{\raggedright\arraybackslash}X>{\raggedright\arraybackslash}X>{\raggedright\arraybackslash}X}
\toprule
\textbf{叙事事件} & \textbf{首次告警时点} & \textbf{两方案判型／预案} & \textbf{结果解释} \\
\midrule
2026 年 3 月 & 事件窗内第 1 日 & 深慢型／深慢预案 & 叙事对照为浅快型，属于判型不匹配，而非漏报 \\
2026 年 7 月 & 实际为起点前 27 日；图中最多展示提前 20 日 & 无法判定／未切换 & 识别到告警状态，但确认与类型证据不足以启动预案 \\
2026 年 10 月 & 起点前 20 日 & 随访不足，不能评价判型 & 数据截至 2026-11-03，不能将截断误记为判别失败或成功 \\
\bottomrule
\end{tabularx}
\end{center}
这三段说明“告警识别”和“类型判别”必须分开验收：前者在重放中全部命中，后者仍分别暴露误判、无法判定和数据尚未覆盖后续走势的问题。验证期已经用于既往评估，本轮重放不构成新的独立样本外证明。由于 R 的重放与修订前逐行逐列一致，本小节的验证期结论对 R 同样成立。\textbf{因此，方案二只保留为重点候选，每段告警重新开放观察期只保留为候选修订，方案一仍是冻结基准；是否替换或采纳，须由新增前瞻数据与另行预登记的判型增强检验决定。}\par
\textbf{（八）C2 否定对照的完整口径。}该“拓扑指标恢复常态后提前解除告警”的候选规则（研发代号 C2）误报日由 778 降至 705，但误报段由 83 增至 94，误报成本由 \(-\)68.4 恶化至 \(-\)114.1，因此净损益差仅为 +30.2。未采用的原因是\textbf{没有净收益增量}，不是误报日增加。该结果仅保留为否定对照。\par
\textbf{（九）拓扑特征的条件判别增量。}在最终产品 S3＋B10 的告警状态下，预登记追加诊断以 272 个真告警早期日（来自 27/45 个事件）与 776 个误报日作主对比（该诊断按原基准日历口径运行，未随本版日历修正与其余候选重算）：11 项指标中 9 项达到预登记条件判别标准，最高受试者工作特征曲线下面积约为 0.744（该指标越接近 1，越能区分真告警早期日与误报日）；5 项在既有七项确认得分之上具有增量判别能力，分别为市场状态平均两两相关、第一特征值占比，以及相关网络二维空腔数量、二维持久景观一范数和行业平均相关。原先“危机邻近日对常态日”的 11 项平均区分度为 0.600，在“真告警早期日对误报日”的告警状态口径下为 0.672；由于两组对照的正负样本定义不同，这只是描述性口径比较，说明旧口径可能低估告警状态下的可分性，不能解释为同一任务上的业绩提升。\textbf{最终产品将拓扑摘要用于告警状态下的信息增强与状态判别，而非独立核心预测模型；新增指标若进入后续版本，仍须另行预登记，不能由该诊断直接改写冻结配置。}\par
\subsection{5.5　已知局限}
当前模型存在四类明确局限，后续研究须分别处理，不作过度承诺：\par
\begin{enumerate}
\item \textbf{前视偏差与损益核算边界。}历史重放核算以事后已知的“事件结束日”为输入，回补起点、止损监测终点、认沽到期与切换资格四个锚点均依赖该信息，因此当前净损益差合计是样本内理论上限，并非实盘可直接实现的收益。核心核算未计一般换手摩擦与滑点，仅计入已建模的仓位变化、止损和认沽现金流。方案一的 +79.0、方案二的 +104.9 及其他对照数字，均为逐事件单位净值损益差与误报段损益差的算术和，不是单一组合累计收益率，也未处理重叠事件的资金占用。五日静默撤销已用于告警状态机，但尚未作为组合解除规则接入损益核算，解除过早或过晚的误差成本仍缺失。
\item \textbf{样本量、窗口与验证边界。}事件级样本仅 45 个，日频样本存在簇内自相关，有效样本量小于名义样本量。45/45 使用“起点前 20 日至事件结束后 20 日”的识别口径，其中包含事件结束后的补充观察期；损益核算的误报归属窗止于事件结束日，核算口径误报与识别窗误报分别为 29.9 和 19.4 日/年，不可混用。主损益核算沿用 6,502 日全历史扫描，其中包含 2025-04-08 以后的日历，不是纯训练期隔离结果。方案一既定规则的封存验证已完成一次；本章方案二补全实验对该区间的再次使用属于描述性重放，不能再视为未被使用过的独立验证集，也不能与主净损益差重复相加。
\item \textbf{类型判别精度与分型证据不足。}三段叙事事件的封存评估中，2026 年 3 月事件出现类型错判，7 月事件无法有效分类，10 月事件因随访不足无法评价；在板块深度下跌而宽基指数小幅波动的结构性行情中，判别力尤其有限。两套事件分型及最终三票方向均使用全样本信息，时间切分只用于早期单特征诊断，不能消除标签与方向中的样本内知识。方案二的补全重放未修复上述问题，因此目前只能保留为重点候选，不能宣称已替代冻结基准或形成可外推的实盘收益。
\item \textbf{候选修订（每段告警重新开放观察期 R）的样本内择优性质。}该修订的提出与评估都使用同一段已知历史：修订方向来自几个因观察期结束而未能切换的窗口（2008 年 1 月、2016 年 1 月、2022 年），改善幅度又在同一批事件上计量。从“根据这些窗口设计补救”到“补救在这些窗口上生效”，全程没有使用新的独立样本，因此属于根据样本内表现选择规则。四点警示：其一，改善高度集中，两个主要配置各 +35.8 中有 33.5 来自 2008 与 2016 两场崩盘开局，事件 12 仅 +1.1；其二，类型修正机制并非普遍有益，更严门槛的配置上同一修正出现负贡献（7 项 5 票组事件 119 修正 \(-\)7.5、7 项 6 票组事件 118 \(-\)3.2），主要配置上的中性偏正含门槛依赖成分；其三，封存重放与修订前逐行逐列一致——2026 年三段确认证据始终不足的窗口在 R 下依然无法切换，验证期又已被方案二补全实验使用过，R 的重放同样不能再视为未被使用过的独立验证集，修订不构成对类型判别短板的修复；其四，方案二＋R 的 +140.7 在两个候选各自择优的基础上叠加，选择偏差亦随之叠加。因此 R 及其全部组合只保留为候选记录：其数字与 +79.0／+104.9 同为历史重放上限口径，证据等级低于冻结基准；正式采纳须有新增前瞻数据与另行预登记的判型增强检验支持。
\end{enumerate}
\section{第六章　系统实现与留出样本验证}
\subsection{6.1　基于危机类型判别的动态风控系统}
第五章所述冻结基准方案采用「触发识别—状态判别—告警管理—绩效核算」四层架构。它的目标不是预测危机起点，而是\textbf{在告警出现后持续检查证据，并在条件满足时切换到对应风控预案}。\par
\textbf{每日收盘后的运行流程}如下：
\begin{enumerate}
\item \textbf{检查是否告警}：最近 5 个交易日简单收益之和不超过 \(-\)4\%、20 日已实现波动率高于滚动 756 日历史的 70 分位且有效历史观测不少于 250 个、20 日回撤达到 5\%，三项至少两项成立；
\item \textbf{在告警状态下继续观察}：逐日计算 7 项状态确认指标与 3 项类型指标，分别判断异常程度与当前市场形态；
\item \textbf{证据充分才调仓}：连续 2 日得到同向类型判断，且第二日确认指标达到门槛时切换预案，最多允许 1 次类型修正；
\item \textbf{管理告警的退出与重启}：连续 5 日无原始命中即撤销，撤销后冷却 10 个交易日，减少重复激活；
\item \textbf{记录损益}：逐日复利记录预案路径，并将危机事件与误报段的损益差分别计入净损益差合计。
\end{enumerate}
\textbf{预案参数表}（预登记后固定）：\par
\begin{center}\small
\begin{tabularx}{\textwidth}{>{\raggedright\arraybackslash}X>{\raggedright\arraybackslash}X>{\raggedright\arraybackslash}X>{\raggedright\arraybackslash}X>{\raggedright\arraybackslash}X}
\toprule
\textbf{危机类型} & \textbf{目标仓位} & \textbf{对冲工具} & \textbf{止损阈值} & \textbf{风控目标} \\
\midrule
深慢型（尺度放大主导） & 50\% & 虚值认沽（执行价 0.95\(\times\)入场价、名义 0.5、权利金 1.5\%） & 20\% & 降低仓位并配置认沽对冲，控制持续回撤风险 \\
浅快型（相关松弛主导） & 80\% & 无 & 12\% & 保持较高仓位并设置相对较紧止损，控制短期冲击 \\
无法判定 & 维持 & 无 & 原规则 & 暂不调整，等待进一步确认 \\
\bottomrule
\end{tabularx}
\end{center}
两档预案的区别在于保护强度：深慢型对应更大的仓位压降，并配置认沽对冲；浅快型保留较高仓位，主要用较紧止损控制短期冲击。\textbf{两个名称沿用第五章的标签约定，实时规则识别的是市场形态，不是预先判断事件最终会跌多深、持续多久。}\par
\textbf{冻结基准方案的样本内结果}（上限口径）：按“事件起点前 20 日至结束后 20 日”的识别口径，45 次危机均被识别，其中 24 次在起点前触发、7 次在事件窗内触发、14 次在事件结束后的补充观察期才触发；31/45 个事件来得及启动预案核算，19/45 个事件实际切换预案。告警日占比为 23.5\%，净损益差合计为 +79.0，其中危机事件内为 +147.4、误报段为 \(-\)68.4。该数值是各事件与各误报段单位净值损益差的算术和，\textbf{不是组合收益率}；无状态管理与候选增强结果只用于第五章对照，不进入本产品的结果汇总。在已知事件起点的机制诊断实验中，正确预案相对无预案的收益差中位数为 +1.82，说明分类正确时预案具有价值，但不代表完整告警管理规则下的逐事件实时收益。\par
\textbf{预留区间的封存验证}（2025-04-08 至 2026-11-03）：六段危机样本均被识别，其中三段按叙事口径认定，三段按固定标签规则重新识别；两段告警提前约 20—23 个交易日。事件识别窗口外的误报约为 3.8 日/年，占交易日的 1.6\%，训练期同口径为 29.9 日/年；2026 年误报为零。但类型判别仍有明显短板：3 月事件误判类型，7 月事件未判定，10 月事件因数据尚未覆盖后续走势而无法评价。\textbf{实盘应用前仍需完成}：另行预登记并优化类型判别规则；将绩效核算中的解除规则改造成可实时执行的规则，消除对事后事件结束日的依赖。\par
\subsection{6.2　系统设计原则}
三项原则贯穿上述系统：\textbf{优先避免漏掉危机}，因为未采取保护的代价可能远高于多防一次的成本；\textbf{控制误报造成的实际调仓成本}，而不是只追求更少的告警日；\textbf{先保证危机识别覆盖，再压低误报}。拓扑特征的工程定位始终是\textbf{信息增强模块}：用于告警后的状态判别，不用于告警撤销，也不作为独立的核心预测模型。研究负责解释机制与适用边界，系统负责执行预案并核算成本；所有绩效数字在通过更严格的样本外检验前，均只代表既定口径下的历史模拟结果，不构成实盘收益承诺。\par
\section{结论}
本文对「拓扑数据分析能否识别并应对 A 股危机」给出三层答案。\par
\textbf{证据层：拓扑摘要没有显示出稳定的独立预测增量。}预登记的增量检验未达到判定标准，增量为 \(-\)0.0076、置信区间跨零；券商公开规则也未通过可交易性检验。这些负面结果只说明相应技术路线没有达到预设要求，不能据此断言拓扑特征不包含任何额外市场信息。\par
\textbf{机制层：部分拓扑响应确实随危机类型而变化。}金融分型（9/36）与拓扑画像分型（15/30）的一致率为 73.3\%；四项拓扑响应通过假发现率校正，并呈现方向相反的类型差异——窗口均方根、平均两两相关和第一特征值占比在深慢型上升、浅快型下降，相关网络零维死亡均值则相反。该结果仅是 45 个事件的样本内证据，不能概括为所有拓扑指标都同向变化，也不构成样本外预测能力证明。\par
\textbf{工程层：冻结基准方案能够把识别、预案与成本核算连成完整流程。}方案一由三项条件取二的触发规则、五日撤销与十日冷却的告警管理、7 项状态确认与 3 项类型投票，以及深慢／浅快两套风控预案组成。按包含事件结束后 20 日补充观察期的识别口径，45 次危机均被覆盖：24 次在起点前触发，7 次在事件窗内触发，14 次在事件结束后才触发；31/45 个事件来得及启动核算，19/45 个事件实际切换预案。告警日占比为 23.5\%，净损益差合计为 +79.0，其中危机事件内为 +147.4、误报段为 \(-\)68.4。该净损益差是逐事件与逐误报段单位净值损益差的算术和，\textbf{不是组合收益率}，并受事件结束日前视偏差与未计一般换手摩擦等限制。\par
候选结果与基准结果必须分开阅读。方案二检验更严格的确认层能否减少误报后的实际调仓，R 检验每段告警重新开放 20 日观察期能否挽回证据晚到的事件。在候选组合中，方案二与分段重新计时规则的组合取得最高样本内净损益差（+140.7）；由于两项修订均经过样本内选择，且验证期未显示类型判别改善，本文将其定位为后续前瞻验证的首选组合，而不替代冻结基准方案。无状态管理、旧触发规则及未采纳配置只用于第五章的消融与选型。后续工作包括优化类型判别规则，以及将绩效核算中的解除规则改造成可实时执行的规则，均须另行预登记并使用新增前瞻数据检验。\par
\section{附录　本文查证的金融侧文献及其结论}
本附录汇总研究过程中查证的金融侧（非 TDA 方法）文献。检索与核对截至 2026-11；文中绩效均为原作者报告，本项目仅对标注「已核对」者做独立复算。\par
\begin{center}\small
\begin{tabularx}{\textwidth}{>{\raggedright\arraybackslash}X>{\raggedright\arraybackslash}X>{\raggedright\arraybackslash}X>{\raggedright\arraybackslash}X}
\toprule
\textbf{文献} & \textbf{主题} & \textbf{核心结论} & \textbf{本文用处} \\
\midrule
Scheffer 等（2009，Nature） & 临界慢化与早期预警 & 系统接近分岔点时恢复力下降，方差、自相关上升，可作通用前兆 & 危机预警的理论纲领 \\
Boettiger 与 Hastings（2012，Proc. R. Soc. B） & 检察官谬误 & 只在已知危机上验证预警指标会严重高估效能 & 全文方法论：预登记、45 个事件的完整样本评估、停止规则 \\
Boettiger 与 Hastings（2012，J. R. Soc. Interface） & 预警检测的极限 & 定量给出预警信号可检测性的边界 & 第三章负面结果的判读边界 \\
Ang 与 Chen（2002，JFE） & 不对称相关 & 熊市中股票间相关性系统性上升，分散化效果下降 & 5.2 节「相关性上升」维度的金融依据 \\
Kritzman 等（2011，J. Portfolio Mgmt.） & 吸收比率 & 第一主成分解释方差占比上升是系统性风险积聚的信号 & 5.2 节首特征值占比指标的金融依据 \\
Billio 等（2012，JFE） & 连通性与系统性风险 & 主成分与格兰杰因果网络可度量金融部门连通性 & 相关结构类指标的文献锚点 \\
Cont（2001，Quantitative Finance） & 收益的典型事实 & 波动聚集、厚尾、相关性时变等风格化事实 & 风险指标基线 M0 的设定依据 \\
Scheffer 等之外：Sornette 等（LPPLS 谱系） & 内生泡沫 & 对数周期幂律奇异性刻画泡沫临近 & 文献池中「内生危机」一支的代表（未复算） \\
\bottomrule
\end{tabularx}
\end{center}
\begin{center}\small
\begin{tabularx}{\textwidth}{>{\raggedright\arraybackslash}X>{\raggedright\arraybackslash}X>{\raggedright\arraybackslash}X>{\raggedright\arraybackslash}X}
\toprule
\textbf{文献} & \textbf{主题} & \textbf{核心结论} & \textbf{本文用处} \\
\midrule
Guritanu 等（2025，Computers） & TDA 危机检测协议 & 严格因果决策规则 + 基线 + 消融；明示「未作交易评估」 & 协议模板；本文补其经济价值空白（已核对其框架） \\
Baitinger 与 Flegel（2021，FMPM） & 拓扑湍流指数 & 相邻持久图距离可预警不利状态；样本外优于其他流行的湍流指数（原作者报告） & 本文复现核对：部分表现属方差尺度效应（已核对） \\
Gidea 与 Katz（2018，Physica A） & 持久景观与崩盘 & 危机前景观范数上升，峰在崩盘中 & 范式来源；第三章增量检验的对立假设来源 \\
兴业证券（2025） & 孔洞数规则 & \(N \le 1\) 做多宣称可观收益（原作者报告） & 第四章核对对象：日度信息系数为正，但未通过可交易性检验（已核对） \\
证监会与市场惯例口径 & 技术性熊市 & 回撤 \(\geq\)20\% 的常见定义；与统计危机标签是两层口径 & 3.1 节标签设计；两层口径不混算 \\
\bottomrule
\end{tabularx}
\end{center}
\textbf{文献池总评}（共 94 条）。检索与核对结果分别见以下两份报告：\par
\noindent\texttt{tda\_systemic\_risk\_literature\_20261030.html}\par
\noindent\texttt{tda\_quant\_finance\_research\_20261028.html}\par
拓扑金融应用中危机检测方向的研究框架最成熟，但严格预警证据不多；最大误区是把「危机检测」当「危机预警」、把「数学稳定性」当「预测稳定性」。证据评级全领域无 A 级。本文的定位由此确定：不做「拓扑预测危机」的第 N 篇正面叙述，做一次同时覆盖增量贡献、危机异质性与经济价值的实证检验和系统实现。\par
\section{参考文献}
\noindent\hangindent=1.6em\hangafter=0 [1] Gidea, M. Topology Data Analysis of Critical Transitions in Financial Networks. NetSci-X, 2017, pp. 47–59.\par
\noindent\hangindent=1.6em\hangafter=0 [2] Gidea, M., Katz, Y. Topological Data Analysis of Financial Time Series: Landscapes of Crashes. Physica A, 2018, 491: 820–834.\par
\noindent\hangindent=1.6em\hangafter=0 [3] Guo, H. et al. Empirical Study of Financial Crises Based on Topological Data Analysis. Physica A, 2020, 558: 124956.\par
\noindent\hangindent=1.6em\hangafter=0 [4] Aguilar, A., Ensor, K. Topology Data Analysis Using Mean Persistence Landscapes in Financial Crashes. Journal of Mathematical Finance, 2020, 10: 648–678.\par
\noindent\hangindent=1.6em\hangafter=0 [5] Ismail, M. S. et al. Early Warning Signals of Financial Crises Using Persistent Homology. Physica A, 2022, 586: 126459.\par
\noindent\hangindent=1.6em\hangafter=0 [6] Baitinger, E., Flegel, S. The Better Turbulence Index? Forecasting Adverse Financial Markets Regimes with Persistent Homology. Financial Markets and Portfolio Management, 2021, 35: 277–308.\par
\noindent\hangindent=1.6em\hangafter=0 [7] Ruiz-Ortiz, M. A. et al. A Persistent-Homology-Based Turbulence Index \& Some Applications of TDA on Financial Markets. arXiv:2203.05603, 2022.\par
\noindent\hangindent=1.6em\hangafter=0 [8] Yen, P. T.-W., Cheong, S. A. Using TDA and Persistent Homology to Analyze the Stock Markets in Singapore and Taiwan. Frontiers in Physics, 2021, 9: 572216.\par
\noindent\hangindent=1.6em\hangafter=0 [9] Yen, P. T.-W., Xia, K., Cheong, S. A. Understanding Changes in the Topology and Geometry of Financial Market Correlations During a Market Crash. Entropy, 2021, 23(9): 1211.\par
\noindent\hangindent=1.6em\hangafter=0 [10] Kulkarni, S. et al. Investigation of Indian Stock Markets Using Topological Data Analysis and Geometry-Inspired Network Measures. Physica A, 2024, 643.\par
\noindent\hangindent=1.6em\hangafter=0 [11] Rai, A. et al. Identifying Extreme Events in the Stock Market: A Topological Data Analysis. Chaos, 2024, 34: 103106.\par
\noindent\hangindent=1.6em\hangafter=0 [12] Valdivia, A. D. Topological Variability in Financial Markets. Quantitative Finance and Economics, 2023, 7(3): 391–402.\par
\noindent\hangindent=1.6em\hangafter=0 [13] Guritanu, E., Barbierato, E., Gatti, A. Topological Machine Learning for Financial Crisis Detection: Early Warning Signals from Persistent Homology. Computers, 2025, 14(10): 408.\par
\noindent\hangindent=1.6em\hangafter=0 [14] Yao, J. et al. Change Point Detection in Financial Market Using Topological Data Analysis. Systems, 2025, 13(10): 875.\par
\noindent\hangindent=1.6em\hangafter=0 [15] Cong, Y. et al. Topological Data Analysis of China's Stock Market Risks to Detect Early Warning Signals. Applied Economics, 2022.\par
\noindent\hangindent=1.6em\hangafter=0 [16] Lamothe-Fernández, P. et al. Topology-Based Machine Learning and Regime Identification in Stochastic, Heavy-Tailed Financial Time Series. Mathematics, 2026, 14(7): 1098.\par
\noindent\hangindent=1.6em\hangafter=0 [17] Katz, Y. A., Biem, A. Time-Resolved Topological Data Analysis of Market Instabilities. Physica A, 2021, 571: 125816.\par
\noindent\hangindent=1.6em\hangafter=0 [18] Akingbade, S. W. et al. Why Topological Data Analysis Detects Financial Bubbles? Communications in Nonlinear Science and Numerical Simulation, 2024.\par
\noindent\hangindent=1.6em\hangafter=0 [19] Souto, H. G. Topological Tail Dependence: Evidence from Forecasting Realized Volatility. Journal of Finance and Data Science, 2023, 9: 100107.\par
\noindent\hangindent=1.6em\hangafter=0 [20] Rudkin, S., Qiu, W., Dłotko, P. Uncertainty, Volatility and the Persistence Norms of Financial Time Series. Expert Systems with Applications, 2023, 223.\par
\noindent\hangindent=1.6em\hangafter=0 [21] Boettiger, C., Hastings, A. Early Warning Signals and the Prosecutor's Fallacy. Proceedings of the Royal Society B, 2012.\par
\noindent\hangindent=1.6em\hangafter=0 [22] Boettiger, C., Hastings, A. Quantifying Limits to Detection of Early Warning for Critical Transitions. Journal of the Royal Society Interface, 2012, 9.\par
\noindent\hangindent=1.6em\hangafter=0 [23] Scheffer, M. et al. Early-Warning Signals for Critical Transitions. Nature, 2009, 461: 53–59.\par
\noindent\hangindent=1.6em\hangafter=0 [24] Ang, A., Chen, J. Asymmetric Correlations of Equity Portfolios. Journal of Financial Economics, 2002, 63(3): 443–494.\par
\noindent\hangindent=1.6em\hangafter=0 [25] Kritzman, M., Li, Y., Page, S., Remington, B. Principal Components as a Measure of Systemic Risk. Journal of Portfolio Management, 2011, 37(2): 112–126.\par
\noindent\hangindent=1.6em\hangafter=0 [26] Billio, M. et al. Econometric Measures of Connectedness and Systemic Risk in the Finance and Insurance Sectors. Journal of Financial Economics, 2012, 104(3).\par
\noindent\hangindent=1.6em\hangafter=0 [27] Cont, R. Empirical Properties of Asset Returns: Stylized Facts and Statistical Issues. Quantitative Finance, 2001, 1: 223–236.\par
\noindent\hangindent=1.6em\hangafter=0 [28] Bubenik, P. Statistical Topological Data Analysis Using Persistence Landscapes. Journal of Machine Learning Research, 2015, 16: 77–102.\par
\noindent\hangindent=1.6em\hangafter=0 [29] Chazal, F., Michel, B. An Introduction to Topological Data Analysis: Fundamental and Practical Aspects for Data Scientists. Frontiers in Artificial Intelligence, 2021.\par
\noindent\hangindent=1.6em\hangafter=0 [30] Wasserman, L. Topological Data Analysis. Annual Review of Statistics and Its Application, 2018, 5: 501–532.\par
\noindent\hangindent=1.6em\hangafter=0 [31] Goel, A., Pasricha, P., Kanniainen, J. Risk Reduced Sparse Index Tracking Portfolio: A Topological Data Analysis Approach. Omega, 2024, 138.\par
\noindent\hangindent=1.6em\hangafter=0 [32] Ismail, M. S. et al. Predicting Next Day Direction of Stock Price Movement Using Machine Learning Methods with Persistent Homology. Applied Soft Computing, 2020, 93: 106422.\par
\noindent\hangindent=1.6em\hangafter=0 [33] Huang, D. et al. Exploring Applications of Topological Data Analysis in Stock Index Movement Prediction. arXiv:2411.13881, 2024.\par
\noindent\hangindent=1.6em\hangafter=0 [34] 郑兆磊, 宫民. 捕捉市场的形状：拓扑数据分析在投资中的应用. 兴业证券金融工程报告, 2025.\par
\noindent\hangindent=1.6em\hangafter=0 [35] Tauzin, G. et al. giotto-tda: A Topological Data Analysis Toolkit for Machine Learning and Data Exploration. Journal of Machine Learning Research, 2021, 22(39): 1–6.\par
\noindent\hangindent=1.6em\hangafter=0 [36] Gokhman, R., Gidea, M. Wavelet-Enhanced Topological Detection of Financial Bubbles. SIAM Conference on Financial Mathematics and Engineering, 2025.\par
\noindent\hangindent=1.6em\hangafter=0 [37] Ismail, M. S. et al. Evidence From Different Correlation Tests: Early Warning Signals of Financial Crises Using Persistent Homology. Frontiers in Applied Mathematics and Statistics, 2022, 8.\par
\noindent\hangindent=1.6em\hangafter=0 [38] Edelsbrunner, H., Harer, J. Computational Topology: An Introduction. American Mathematical Society, 2010.\par
\noindent\hangindent=1.6em\hangafter=0 [39] Cohen-Steiner, D., Edelsbrunner, H., Harer, J. Stability of Persistence Diagrams. Discrete \& Computational Geometry, 2007, 37: 103–120.\par
\noindent\hangindent=1.6em\hangafter=0 [40] Carlsson, G. Topology and Data. Bulletin of the American Mathematical Society, 2009, 46: 255–308.\par
\textbf{研究纪律}：先预登记后计算 · 独立重复运行一致性校验 · 不根据结果反向调整参数 · 不事后修改判定 · 如实报告阴性结果。\par
\textbf{数据与代码}：工程目录 tda\_project/；预登记文件存于 config/；实验输出与报告存于 outputs/。训练期及全历史重放结果为样本内描述性统计与历史重放损益评估；方案一验证期（2025-04-08 起）已完成一次封存样本外检验，后续方案二补全实验对该区间的重用仅作描述性重放，不能再视为未被使用过的独立验证集。历史重放损益核算以事件结束日为已知输入，存在前视偏差，相关损益数值仅为样本内上限。\par
\end{document}
